% Master's Thesis — José Gonçalves Lima Neto
% IME-USP, Computer Science, 2026
%
% Based on the IME-USP LaTeX template (ccsl-usp/modelo-latex).
% Template authors: Jesús P. Mena-Chalco, Fabio Kon, Paulo Feofiloff, Nelson Lago.
% Template license: MIT. Thesis content: CC-BY 4.0.


%%%%%%%%%%%%%%%%%%%%%%%%%%%%%%%%%%%%%%%%%%%%%%%%%%%%%%%%%%%%%%%%%%%%%%%%%%%%%%%%
%%%%%%%%%%%%%%%%%%%%%%%%%%% DOCUMENT CLASS & PACKAGES %%%%%%%%%%%%%%%%%%%%%%%%%%
%%%%%%%%%%%%%%%%%%%%%%%%%%%%%%%%%%%%%%%%%%%%%%%%%%%%%%%%%%%%%%%%%%%%%%%%%%%%%%%%

% The LAST language listed becomes the document default.
% Swapped order: English is now the primary language.
% Brazilian Portuguese is still loaded (required by IME-USP for the resumo).
\documentclass[a4paper,12pt,twoside,brazilian,english]{book}

\usepackage{imegoodies}
\usepackage[thesis]{imelooks}

% Directories where figures are located
\graphicspath{{figures/},{template/logos/}}

% Framed placeholders for screenshots until final figures are exported
\newcommand{\placeholder}[2][0.85\textwidth]{%
  \fbox{\parbox{#1}{\centering\vspace{2cm}[SCREENSHOT PLACEHOLDER: #2]\vspace{2cm}}}}

% Quick language-switching commands (\en, \br, \texten{}, \textbr{})
\babeltags{br = brazilian, en = english}


%%%%%%%%%%%%%%%%%%%%%%%%%%%%%%%%%%%%%%%%%%%%%%%%%%%%%%%%%%%%%%%%%%%%%%%%%%%%%%%%
%%%%%%%%%%%%%%%%%%%%%%%%%%%%%%%%%%% METADATA %%%%%%%%%%%%%%%%%%%%%%%%%%%%%%%%%%%
%%%%%%%%%%%%%%%%%%%%%%%%%%%%%%%%%%%%%%%%%%%%%%%%%%%%%%%%%%%%%%%%%%%%%%%%%%%%%%%%

\addbibresource{references.bib}

\babelhyphenation{documentclass latexmk soft-ware clsguide}
\babelhyphenation[brazilian]{hi-pó-te-ses ar-qui-te-tu-rais}
\babelhyphenation[english]{Arch-Hypo Hypo-Stage}

% ---- Title (English primary; translated title is the Portuguese version) ----
\title{Hypo-Stage: Tool-Supported Architectural Hypothesis Engineering\\
in Professional Software Development Teams --- A Case Study}[]
\translatedtitle{Hypo-Stage: Engenharia de Hipóteses Arquiteturais Apoiada por\\
Ferramenta em Equipes Profissionais de Desenvolvimento de Software ---\\
Um Estudo de Caso}[]

\author{José Gonçalves Lima Neto}

\def\profa{Prof\kern.02em.\kern-.07emª\kern.07em}
\def\dra{Dr\kern-.04em.\kern-.11emª\kern.07em}

\orientador{Prof. Dr. Paulo Roberto Miranda Meirelles}

% Confirmed examination committee (titulares). The full roster of titulares
% and suplentes, including institutional affiliations, lives in
% front-matter/banca.tex (single source of truth for defense scheduling).
\banca{
  Prof. Dr. Paulo Roberto Miranda Meirelles (Presidente) -- IME-USP,
  Prof. Dr. Jorge Melegati -- Free University of Bozen-Bolzano (UniBZ),
  Prof. Dr. Marco Túlio de Oliveira Valente -- UFMG,
}

\tipotese{
  mestrado,
  %definitiva, % TODO: uncomment after defense to switch to "final version" title page
  programa={Computer Science},
}

\defesa{
  local={São Paulo},
  data=2026-05-15, % TODO: confirm exact defense date (expected May 2026)
}

% TODO: uncomment and fill when funding source is confirmed
%\apoio{This study was financed in part by the Coordenação de
%      Aperfeiçoamento\\ de Pessoal de Nível Superior -- Brasil
%      (CAPES) -- Finance Code 001}
%
%\apoio{During the development of this work, the author received\\
%      financial support from FAPESP -- grant \#aaaa/nnnnn-d}

\direitos{CC-BY}

% Ficha catalográfica generated at https://fc.ime.usp.br/
% Depends on packages calc, setspace, and ragged2e (all already loaded by imegoodies/imelooks).
\fichacatalografica{%
\begingroup\centering\singlespacing\small
\hyphenrules{nohyphenation}\hbadness=10000
Ficha catalográfica elaborada com dados inseridos pelo(a) autor(a)\\
Biblioteca Carlos Benjamin de Lyra\\
Instituto de Matemática, Estatística e Ciência da Computação\\
Universidade de São Paulo\par
\vspace{2\baselineskip}\hrule\vspace{.8\baselineskip}
\RaggedRightRightskip 0pt plus 30pt minus 0pt\relax
\RaggedRightParfillskip 20pt plus 40pt minus 10pt\relax
\ttfamily\hspace{2em}\begin{minipage}[t]{125mm}
\RaggedRight\sloppy\setlength{\parindent}{\widthof{123}}

\noindent Lima Neto, José Gonçalves

Hypo-Stage: Tool-Supported Architectural Hypothesis Engineering in Professional Software Development Teams --- A Case Study / José Gonçalves Lima Neto; orientador, Paulo Roberto Miranda Meirelles. --- São Paulo, 2026.

87 p.

\vspace{1\baselineskip}


\vspace{1\baselineskip}


Dissertação (Mestrado) --- Programa de Pós-Graduação em Ciência da Computação / Instituto de Matemática, Estatística e Ciência da Computação / Universidade de São Paulo.

Bibliografia

Versão original

\vspace{1\baselineskip}


\vspace{1\baselineskip}


1. Software Architecture. 2. Architectural Uncertainty. 3. Hypotheses Engineering. 4. ArchHypo. 5. Hypo-Stage. 6. Internal Developer Portal. 7. Backstage. 8. Case Study. 9. Empirical Evaluation. I. Miranda Meirelles, Paulo Roberto. II. Prof. Dr.

\end{minipage}\par
\vspace{1\baselineskip}\hrule\vspace{.5\baselineskip}\rmfamily
Bibliotecárias do Serviço de Informação e Biblioteca\\
Carlos Benjamin de Lyra do IME-USP, responsáveis pela\\
estrutura de catalogação da publicação de acordo com a AACR2:\\
Maria Lúcia Ribeiro CRB-8/2766; Stela do Nascimento Madruga CRB 8/7534.
\par\endgroup
}

%\nocolorlinks % TODO: uncomment for black-and-white printing

% Bump numbering and TOC depth to 3 so \subsubsection entries (used for the
% emerging-pattern catalogue in Chapter~5) are numbered and listed in the
% summary. The imelooks template defaults both counters to 2.
\setcounter{secnumdepth}{3}
\setcounter{tocdepth}{3}


%%%%%%%%%%%%%%%%%%%%%%%%%%%%%%%%%%%%%%%%%%%%%%%%%%%%%%%%%%%%%%%%%%%%%%%%%%%%%%%%
%%%%%%%%%%%%%%%%%%%%%%%%%%%%%%%% DOCUMENT BODY %%%%%%%%%%%%%%%%%%%%%%%%%%%%%%%%%
%%%%%%%%%%%%%%%%%%%%%%%%%%%%%%%%%%%%%%%%%%%%%%%%%%%%%%%%%%%%%%%%%%%%%%%%%%%%%%%%

\begin{document}

%%%%%%%%%%%%%%%%%%%%%%%%%%%%% FRONT MATTER %%%%%%%%%%%%%%%%%%%%%%%%%%%%%%%%%%%%%

\frontmatter
\pagestyle{plain}
\onehalfspacing

\maketitle

%%% Advisor-signed original-version declaration (pre-deposit) %%%
% !TeX root = ../thesis.tex
% Advisor-signed original-version declaration (pre-deposit).
% Appears after \maketitle (IME-USP cover + title + ficha verso) and before
% the dedicatória. The text mirrors the "original version" declaration
% already rendered on the IME-USP title page by \@versionInfoBlock, and
% adds the advisor signature block required for deposit approval.

\cleardoublepage
\thispagestyle{empty}
\vspace*{\fill}

\noindent This is the original version of the thesis prepared by candidate
José Gonçalves Lima Neto, as submitted to the Examining Committee.

\vspace{3cm}

\begin{center}
    \rule{8cm}{0.4pt}\\[2pt]
    Prof. Dr. Paulo Roberto Miranda Meirelles\\
    \textit{Advisor}\\
    Institute of Mathematics, Statistics, and Computer Science --- University of São Paulo
\end{center}

\vspace{1cm}

\begin{center}
    São Paulo, \rule{1.5cm}{0.4pt} of \rule{3.5cm}{0.4pt} of 2026.
\end{center}

\vspace*{\fill}
\clearpage


%%% Dedication %%%
% !TeX root = ../thesis.tex
% Dedicatória — rendered via the IME-USP `dedicatoria` environment, which
% handles page geometry, typography, and the forced-recto placement.

\begin{dedicatoria}
To my wife and family,\\
whose unwavering support made this journey possible.
\end{dedicatoria}


% Note: \frontmatter (above) already enables roman page numbering, so no
% explicit \pagenumbering{roman} is required here. Adding one would reset
% the page counter and create duplicate hyperref anchors (name{page.v}).

%%% Acknowledgments %%%
\chapter*{Acknowledgments}

\begin{flushright}
\itshape
Good decisions come from experience.\\
Experience comes from making bad decisions.\\
\upshape --- Mark Twain
\end{flushright}

\vspace{1.5em}

Special thanks to my advisor, Prof.\ Dr.\ Paulo Meirelles, for his exceptional guidance and
support throughout this work. His mentorship shaped not only the direction of this research
but also the way in which it was conducted. I also extend my gratitude to the staff of
\mbox{IME-USP} for their assistance and institutional support, and to the teams who
participated in this study for their openness, generosity with their time, and genuine
engagement with the research. Finally, I thank my wife and family for their unwavering
motivation and encouragement --- without them, none of this would have been possible.

%%% Abstract (EN primary, PT secondary — both required by IME-USP) %%%
%!TeX root=../thesis.tex

% Keywords must be defined before the abstract/resumo commands.
\keywords{Software Architecture, Architectural Uncertainty, Hypotheses Engineering, ArchHypo, Hypo-Stage, Internal Developer Portal, Backstage, Case Study, Empirical Evaluation}

\palavraschave{Arquitetura de Software, Incerteza Arquitetural, Engenharia de Hipóteses, ArchHypo, Hypo-Stage, Portal do Desenvolvedor Interno, Backstage, Estudo de Caso, Avaliação Empírica}

% English abstract (primary)
\abstract{
Software engineering teams in continuous-delivery organizations routinely face
architectural decisions under incomplete information, yet few lightweight tools
exist to help them record, prioritize, and learn from those uncertainties
collaboratively. The ArchHypo technique applies hypotheses engineering to
software architecture: teams formulate architectural uncertainties as
falsifiable hypotheses, assess each by uncertainty level and impact, and
define a technical plan for handling it. Prior evaluations demonstrated the
technique's value but reported a steep learning curve and dependence on expert
facilitation, both exacerbated by the absence of tool support.

This dissertation presents the first empirical account of a tool-mediated
application of ArchHypo in a professional software development organization. The artifact is Hypo-Stage,
an open-source Backstage plugin deployed on a large organization's Internal
Developer Portal, where engineers already work. The study adopts a flexible,
multiple case-study design with formative evaluation and action-research
elements, following the framework of \citet{robson2016}. Two teams --- a
Product team of twelve engineers and a Platform team of four --- used Hypo-Stage
over two agile sprints totalling four weeks. Data were collected through the
author's participant observation (continuous field notes), artifact analysis of
the thirteen hypotheses and associated technical plans registered in Hypo-Stage,
and open-ended active listening sessions with the senior technical leaders of
each team. Analysis followed the three-stage thematic coding procedure
described by \citet{robson2016} for qualitative data in flexible designs.

Four contributions emerge from the work.
\textbf{(1)} A detailed observational account of how two teams in a large-scale
product organization adopted and used Hypo-Stage, characterizing hypothesis
registration dynamics, role stratification, and tool adoption patterns.
\textbf{(2)} A characterization of the architectural uncertainties the teams
made explicit: their sources, quality attributes (dominated by Reliability and
Auditability), uncertainty-impact profiles, and technical plan strategies.
\textbf{(3)} Five novel technical architectural hypothesis archetypes ---
\textit{Hypothesis Cluster}, \textit{Regulatory Trigger Hypothesis},
\textit{Decoupling Seam Hypothesis}, \textit{Load-Validated Hypothesis}, and
\textit{Shared Resource Governance Hypothesis} --- derived from thematic analysis
of the hypothesis dataset. These are the first empirically grounded
characterization of recurring hypothesis shapes in the ArchHypo research line,
directly addressing a gap identified in prior work; they are presented as
emerging patterns requiring further validation.
\textbf{(4)} Nine categories of empirically grounded refinements to Hypo-Stage,
each traced to a concrete observation and linked to a code commit, including a
hypothesis evolution chart, a prioritization dashboard, and inline formulation
guidance that reduces the technique's learning barrier.
}

% Portuguese abstract (required by IME-USP regulations)
\resumo{
Equipes de engenharia de software em organizações de entrega contínua tomam
rotineiramente decisões arquiteturais sob informação incompleta, mas dispõem de
poucas ferramentas leves para registrar, priorizar e aprender com essas
incertezas de forma colaborativa. A técnica ArchHypo aplica a engenharia de
hipóteses à arquitetura de software: as equipes formulam incertezas
arquiteturais como hipóteses falsificáveis, avaliam cada uma por nível de
incerteza e impacto, e definem um plano técnico para tratá-las. Avaliações
anteriores demonstraram o valor da técnica, mas reportaram uma curva de
aprendizado acentuada e dependência de facilitação especializada, ambos
agravados pela ausência de suporte ferramental.

Esta dissertação apresenta o primeiro relato empírico de uma aplicação mediada
por ferramenta do ArchHypo em uma organização profissional de desenvolvimento de software. O artefato é o Hypo-Stage,
um plugin \textit{open source} para o Backstage, implantado no Portal do
Desenvolvedor Interno de uma grande organização de software, onde os engenheiros
já atuam cotidianamente. O estudo adota um desenho flexível de múltiplos
estudos de caso, com propósito avaliativo formativo e elementos de
pesquisa-ação, seguindo o arcabouço de \citet{robson2016}. Duas equipes --- uma
equipe de Produto com doze engenheiros e uma equipe de Plataforma com quatro ---
utilizaram o Hypo-Stage ao longo de duas \textit{sprints} ágeis, totalizando
quatro semanas. Os dados foram coletados por meio de observação participante
do próprio autor (notas de campo contínuas), análise de artefatos das treze
hipóteses e planos técnicos registrados no Hypo-Stage, e sessões abertas de
escuta ativa com os líderes técnicos seniores de cada equipe. A análise
seguiu o procedimento de codificação temática em três estágios descrito por
\citet{robson2016} para dados qualitativos em designs flexíveis.

Quatro contribuições emergem do trabalho.
\textbf{(1)} Um relato observacional detalhado de como duas equipes em uma
organização de produto em escala adotaram e utilizaram o Hypo-Stage,
caracterizando a dinâmica de registro de hipóteses, a estratificação por papel
e os padrões de adoção da ferramenta.
\textbf{(2)} Uma caracterização das incertezas arquiteturais tornadas explícitas
pelas equipes: suas fontes, atributos de qualidade (dominados por
Confiabilidade e Auditabilidade), perfis de incerteza-impacto e estratégias
de plano técnico.
\textbf{(3)} Cinco arquétipos técnicos de hipóteses arquiteturais ---
\textit{Hypothesis Cluster}, \textit{Regulatory Trigger Hypothesis},
\textit{Decoupling Seam Hypothesis}, \textit{Load-Validated Hypothesis} e
\textit{Shared Resource Governance Hypothesis} --- derivados da análise temática
do conjunto de hipóteses. Estes constituem a primeira caracterização empírica de
formas recorrentes de hipóteses na linha de pesquisa do ArchHypo, abordando
diretamente uma lacuna identificada em trabalhos anteriores; são apresentados
como padrões emergentes que requerem validação adicional.
\textbf{(4)} Nove categorias de refinamentos empiricamente fundamentados ao
Hypo-Stage, cada um rastreado a uma observação concreta e vinculado a um
\textit{commit} no repositório, incluindo um gráfico de evolução das hipóteses,
um painel de priorização e orientações de formulação \textit{inline} que
reduzem a barreira de aprendizado da técnica.
}



%%%%%%%%%%%%%%%%%%%%%%%%%%%%% LISTS & TOC %%%%%%%%%%%%%%%%%%%%%%%%%%%%%%%%%%%%%%

\cleardoublepage

\newcommand\disablenewpage[1]{{\let\clearpage\par\let\cleardoublepage\par #1}}

\bgroup
\raggedbottom

%%% Abbreviations %%%
\disablenewpage{\chapter*{List of Abbreviations}}

\begin{tabular}{rl}
   SA       & Software Architecture\\
   ArchHypo & Architecture Hypothesis (technique)\\
   IDP      & Internal Developer Portal\\
   NFR      & Non-Functional Requirements\\
   QA       & Quality Attribute\\
   DSR      & Design Science Research\\
   HE       & Hypotheses Engineering\\
   IME      & Institute of Mathematics and Statistics (Universidade de São Paulo)\\
   USP      & University of São Paulo\\
   ADR      & Architecture Decision Record\\
   % TODO: add more abbreviations as they appear in the text
\end{tabular}

\clearpage

%%% Auto-generated lists %%%
\disablenewpage{\listoffigures}
\disablenewpage{\listoftables}
%%% List of software programs (repositories central to this thesis) %%%
\disablenewpage{\chapter*{\programlistname}}

\phantomsection
\addcontentsline{toc}{chapter}{\programlistname}

\noindent\textbf{Hypo-Stage} --- open-source Backstage plugin that operationalizes the ArchHypo technique for hypothesis management; empirical artifact of this thesis. Repository: \url{https://github.com/ArchHypo/hypo-stage}.

\par\vspace{\baselineskip}
\noindent\textbf{Backstage} --- open-source framework for Internal Developer Portals (CNCF); organizational platform in which Hypo-Stage is deployed in this study. Repository: \url{https://github.com/backstage/backstage}.

%%% Table of Contents (required) %%%
\tableofcontents

\egroup


%%%%%%%%%%%%%%%%%%%%%%%%%%%%%%%% CHAPTERS %%%%%%%%%%%%%%%%%%%%%%%%%%%%%%%%%%%%%%

\mainmatter
\pagestyle{mainmatter}
\singlespacing

%!TeX root=../thesis.tex

\chapter{Introduction}
\label{ch:intro}
\label{chap:introduction}

Software architecture (SA) plays a central role in the development and evolution
of software-intensive systems. The SA concept structures a system into elements and relations that form the basis of
its organization, conveying coherence, scalability, and maintainability for the
evolving software over time \citep{perryWolf1992}. Despite the ever-changing practices of modern
software development, SA still provides the basis for analyzing the quality
attributes of modern systems, guiding technical decisions made throughout
the development process via a diverse set of architectural methodologies \citep{souza2019}. In organizations that rely on
software as a core part of their business, architecture is not a one-time design
artifact, but rather a continuous decision-making process that must accommodate changing
requirements, evolving technologies, and critical organizational constraints
\citep{silva2024tse}.

In recent years, agile and continuous development practices have intensified this fast-paced pressure. 
Teams are expected to deliver more features in shorter cycles while preserving scalability, performance, and other quality attributes of their systems that continue to grow in size and complexity, as observed in \citep{silva2024tse,silva2024ieeesw}. 
In such environments, architecture decisions are frequently made under uncertainty: solutions or requirements may be incomplete or unstable, the impact of technology choices may not be fully understood, and the behavior of the system under realistic workloads may only become visible after deployment \citep{silva2024tse}. 
Despite such challenges, architectural uncertainties often remain implicit and unmanaged, and decisions tend to rely on the tacit knowledge and intuition of a small group of experienced architects, as summarized in systematic evidence on architectural derivation methods \citep{souza2019}.

At the same time, many organizations seek to organize themselves around autonomous, cross-functional teams responsible for end-to-end delivery of features. 
Recent industry reports, such as the 2024 InfoQ Software Architecture and Design Trends Report \citep{infoq2024trends}, highlight socio-technical architecture as an emerging trend that emphasizes the co-design of technical and organizational structures, and the importance of considering how team boundaries and communication patterns shape system architectures. 
In such a scenario, architecture work is distributed across teams rather than centralized in a separate function or role. Moreover, in these scenarios, the Conway's Law is treated as a design constraint to be intentionally leveraged rather than an accident to be ignored \citep{infoq2024trends}. 
However, when architecture decisions depend on implicit assumptions held by a few experts in the organization, this distribution becomes difficult in practice. Less experienced team members may struggle to participate meaningfully in architectural discussions, and teams may find it hard to reason about the risks and trade-offs of alternative solutions even when they are applying any one of the well-known and validated architectural methodologies mapped in prior work \citep{souza2019}.
This challenge reveals the case of the architectural uncertainty management, which still is a critical aspect of the architecture work, but is often neglected or poorly supported in practice by the available tools and methodologies.

\section{Architectural Uncertainty and Hypotheses Engineering}

Uncertainty in architecture decisions proceeds from limited knowledge about requirements, future usage scenarios, and the properties of candidate solutions as reviewed in \citep{silva2024tse}. 
Traditional risk management techniques assume that relevant events and their probabilities are known or can be estimated, which is not always the case when dealing with evolving systems and emerging technologies \citep{silva2024tse}. 
A complementary perspective is to acknowledge uncertainty explicitly and as an inevitable aspect of the architecture work, and therefore, to manage it as a first-class concern during the architecture work as proposed in \citep{silva2024tse}.

In this regard, Hypotheses Engineering has emerged as an approach to manage uncertainty in software development by treating assumptions as explicit, falsifiable hypotheses that can be prioritized and tested through experiments \citep{silva2024tse,silva2024ieeesw}. In this context, ArchHypo is the technique that adapts hypotheses engineering to software architecture. 
It proposes that teams identify \emph{architectural hypotheses} --- statements capturing uncertainties that may affect architectural decisions --- assess each hypothesis by its uncertainty level and impact, and define a technical plan for handling it \citep{silva2024tse}.

The technical plan may include actions to reduce uncertainty, such as experiments, architectural spikes, or tracer bullets, and actions to reduce the impact of potential changes, such as modularization or flexible deployment strategies \citep{silva2024tse}. 
By iteratively updating hypotheses and technical plans over time, ArchHypo aims to support a more deliberate and distributed approach to architectural decision-making, allowing teams to postpone irreversible decisions when justified and to focus their efforts on the most critical uncertainties \citep{silva2024tse,silva2024ieeesw}.

Empirical studies have started to explore how ArchHypo works in practice. An experience report in a software start-up describes how the technique helped the founder identify and prioritize architectural uncertainties related to availability, usability, and scalability over a ten-month period, guiding the evolution of the product and avoiding premature investments in infrastructure \citep{silva2024ieeesw}. 
A design science study in a mission-critical public-sector system shows that integrating ArchHypo into the development process enabled the team to distribute architectural work across iterations and to make better-informed decisions in the face of changing requirements \citep{silva2024tse}. 
At the same time, both studies report that the technique introduces a non-trivial learning curve and requires adjustments to existing processes, which might pose a challenge to its adoption by the team as reported in \citep{silva2024tse,silva2024ieeesw}.

\section{From Requirements to Architecture and the Role of Tool Support}

The challenges of managing architectural uncertainty are closely related to the long-standing problem of deriving architectural models from requirements specifications. 
\citet{souza2019} reviewed 39 methods for this derivation and concluded that existing approaches still rely heavily on the tacit knowledge of architects, provide limited support for less experienced practitioners, and often lack explicit evaluation mechanisms. 
Many methods describe high-level processes or frameworks, but offer little guidance on how to operationalize them in everyday practice as reported in \citep{souza2019}.

The same study highlighted gaps in tool support. A significant portion of the methods did not include any supporting tool, and when tools were present, they tended to address only parts of the process as also reported in \citep{souza2019}. 
The authors called for more automation, better mechanisms to capture and reuse architectural knowledge, and empirical validation of tools in realistic settings. 
These findings resonate with the observations from ArchHypo studies: while the technique provides a structured way to make architectural uncertainties explicit and to plan actions, applying it in practice is highly demanding without adequate supporting tool \citep{silva2024tse,silva2024ieeesw}.

Grounded in these studies' contributions, \citet{correa2025} developed the Hypo-Stage, a support tool designed to help teams apply the ArchHypo technique in practice.
Up to a specific point in its history, the work focused on designing and implementing the core functionality for enabling ArchHypo capabilities, such as representing architectural hypotheses, recording their uncertainty and impact assessments, and managing associated technical plans. \footnote{The first version of Hypo-Stage and the corresponding source code history up to commit \texttt{ceee509776508081ebdd473c2c4f710b8ef55947} document the initial design and implementation of Hypo-Stage \citep{correa2025}.} 
The work also discussed initial design decisions and limitations, including usability aspects and open questions about how the tool would fit into different team workflows.
Chapter~\ref{chap:background} presents a more detailed description of Hypo-Stage and its associated contributions \citep{correa2025}.

\section{Problem Statement}

When reviewed in intersection, the reviewed studies in this present work outline a clear problem. 
On the one hand, modern software organizations need to manage architectural uncertainty in a way that supports continuous evolution and involves the broader development team. 
On the other hand, empirical evidence and systematic mapping studies indicate that:

\begin{itemize}
  \item Methods for deriving and evolving architectures from requirements still depend strongly on architects' tacit knowledge and lack comprehensive tool support as reported in \citep{souza2019}.
  \item ArchHypo offers a promising structure for handling architectural uncertainty, but teams perceive the technique as hard to learn and requiring process changes, especially without specialized support as reported in \citep{silva2024tse,silva2024ieeesw}.
  \item Hypo-Stage provides an initial implementation of tool support for ArchHypo, but its impact on real software engineering teams and the refinements needed for effective use in industry remain unknown as reported in \citep{correa2025}.
\end{itemize}

In this context, there is a need to understand how software engineering teams,
with diverse roles, tenures, and responsibilities, interact with a tool like
Hypo-Stage when working on their products and services. By addressing this problem, this work seeks to gain an understanding of what architectural uncertainties
teams manage through such the tool, and which aspects of Hypo-Stage and its integration into the development process still require refinement.
Investigating these concerns can contribute both to the practical adoption of ArchHypo via Hypo-Stage and to the broader discussion on how to support teams in sharing architectural responsibilities for accelerating adoption of the socio-technical architecture paradigm mapped by \citep{infoq2024trends}.

\section{Objectives}

The overall objective of this work is to empirically evaluate and refine Hypo-Stage, the tool-supported implementation of the ArchHypo technique introduced in prior Hypo-Stage work \citep{correa2025}, by applying it with diverse professional software development teams in a real-world organizational context that maintains a scaled software product in Latin America. 
More specifically, this work aims to:

\begin{itemize}
  \item Investigate how software engineering teams appropriate and use Hypo-Stage and ArchHypo in their day-to-day activities when working on an existing, scaled product.
  \item Characterize the architectural hypotheses that teams register and manage through Hypo-Stage, identifying emerging patterns among them.
  \item Identify benefits and challenges perceived by team members regarding the use of Hypo-Stage, including their engagement with the tool, usability, and fit with existing practices.
  \item Derive and implement refinements to Hypo-Stage and its usage guidelines based on empirical findings, with the goal of improving its suitability for supporting architectural work in teams.
\end{itemize}

These objectives deliberately build on, rather than replicate, the previous studies' contributions in the ArchHypo research area. 
While \citet{correa2025} focused on designing and implementing Hypo-Stage, this work treats the tool as an existing artifact and concentrates on its evaluation and evolution in a professional software development context.

\section{Research Questions}

To achieve the above objectives, this work is guided by the following research questions, listed in the same order in which they are organized in Chapter~\ref{chap:results}: tool use and appropriation (RQ1); recorded uncertainties and structural patterns in the hypothesis data (RQ2); perceived benefits and challenges (RQ3); and empirically grounded refinements to Hypo-Stage and its usage guidelines (RQ4).

\begin{description}
  \item[RQ1:] How do software engineering teams use Hypo-Stage and the ArchHypo technique when working on a scaled software product in a real-world organization?
  \item[RQ2:] What types and recurring structural shapes of architectural uncertainties do these teams identify, record, and manage using Hypo-Stage, and what patterns emerge from their content and organizational context?
  \item[RQ3:] What benefits and challenges do team members perceive in using Hypo-Stage to support architectural decision-making and uncertainty management?
  \item[RQ4:] Which refinements to Hypo-Stage and its usage guidelines emerge from the empirical study, and how do they address the identified challenges?
\end{description}

Chapter~\ref{chap:research-design} operationalizes these questions through aligned data sources: participant observation and Hypo-Stage artifacts support RQ1 (tool use) and RQ2 (recorded uncertainties); active listening foregrounds RQ3 (perceived benefits and challenges) while also surfacing needs that motivate refinements under RQ4.

\begin{table}[htbp]
\centering
\caption{Alignment between research questions and study objectives.}
\label{tab:rq-objectives}
\begin{tabular}{p{0.08\linewidth} p{0.42\linewidth} p{0.42\linewidth}}
\toprule
\textbf{RQ} & \textbf{Research Question} & \textbf{Primary Objective Addressed} \\
\midrule
RQ1 & How do software engineering teams use Hypo-Stage and ArchHypo when working on a scaled product? & Investigate appropriation and use patterns of Hypo-Stage in day-to-day team activities. \\
\midrule
RQ2 & What types and recurring structural shapes of architectural uncertainties do teams identify, record, and manage using Hypo-Stage, and what patterns emerge from their content and organizational context? & Characterize architectural uncertainties by source, quality attribute, and lifecycle; identify recurring structural archetypes of hypotheses that emerge from the empirical data as a novel contribution to the ArchHypo pattern literature. \\
\midrule
RQ3 & What benefits and challenges do team members perceive in using Hypo-Stage? & Identify perceived benefits and challenges related to engagement with the tool, usability, and workflow fit. \\
\midrule
RQ4 & Which refinements emerge from the empirical study and how do they address identified challenges? & Derive and implement empirically grounded refinements to Hypo-Stage and its usage guidelines. \\
\bottomrule
\end{tabular}
\end{table}

\section{Research Contributions}

By answering its set of research questions, this work attempts to make the following contributions:

\begin{itemize}
  \item An empirical account of how multiple professional software development teams use Hypo-Stage --- the tool-supported implementation of ArchHypo \citep{correa2025} --- while working on a scaled software product, including observed adoption patterns, practices, and obstacles.
  \item A characterization of the architectural uncertainties these teams made explicit and managed through Hypo-Stage, describing their sources, quality attributes, structural archetypes, and evolution.
  \item A set of novel \textit{technical architectural hypothesis patterns} derived from thematic analysis of the hypotheses registered during the study. These patterns --- Hypothesis Cluster, Regulatory Trigger Hypothesis, Decoupling Seam Hypothesis, Load-Validated Hypothesis, and Shared Resource Governance Hypothesis --- constitute the first empirically grounded answer to the question of what recurring shapes architectural hypotheses take when ArchHypo is applied in practice, addressing a gap explicitly identified as future work in the prior ArchHypo pattern literature \citep{silva2024patterns}.
  \item A set of empirically grounded refinements to the Hypo-Stage tool and its usage guidelines, aimed at improveing the tool's usability and engagement, and better integrating the technique into team workflows as reported in \citep{correa2025}.
  \item Reflections on the implications of these findings for supporting distributed architectural work in agile and continuous development settings, connecting the ArchHypo experience to broader concerns about tacit knowledge and tool support in software architecture as reported in \citep{souza2019,infoq2024trends}.
\end{itemize}

\paragraph{Preview of key findings.}
This study produced three families of concrete results.
First, hypothesis registration proved consistently senior-led in both teams:
twelve of the thirteen hypotheses were authored by senior or staff engineers,
and sprint delivery pressure was the dominant barrier to broader participation ---
a finding that directly motivates both the tool improvements of RQ4 and the
organizational pattern A3 (Section~\ref{sec:group-a-patterns}).
Second, the thirteen hypothesis profiles reveal a strong concentration of
compliance-driven and shared-infrastructure uncertainties, producing two of
the five emerging patterns (\textit{Regulatory Trigger Hypothesis} and
\textit{Shared Resource Governance Hypothesis}) that are novel contributions
to the ArchHypo pattern literature.
Third, the refinement cycle produced nine categories of empirically grounded
improvements to Hypo-Stage, including an evolution chart for tracking
assessment changes over time. These findings are previewed here to orient the reader. A complete result analysis and
discussion appear in Chapter~\ref{chap:results} and Chapter~\ref{chap:conclusion}.

\section{Research Structure}

The remainder of this work is organized as follows.

Chapter~\ref{chap:background} provides the theoretical foundations for the study.
It introduces key concepts in software architecture and architectural uncertainty;
presents the ArchHypo technique in detail, including its core concepts, assessment
model, technical plan strategies, and lifecycle; reviews the two prior empirical
evaluations of ArchHypo (at Catch Solve and PropSEP/Jetsoft); describes Internal
Developer Portals and the Backstage platform as the delivery vehicle for Hypo-Stage;
presents the initial Hypo-Stage tool \citep{correa2025}; and surveys systematic evidence on the gap between architectural
derivation methods and tool support \citep{souza2019}.

Chapter~\ref{chap:related-work} synthesizes related literature and industry context relevant to tool-supported architectural decision-making and positions this work within that space.

Chapter~\ref{chap:research-design} presents the research design and methods adopted
in this work. It describes the research strategy (a flexible multiple case study
with formative evaluation and action research elements), the conceptual research
framework that organizes the key constructs of the study, the organizational context and
case selection rationale for the two participating teams (a Product Team and a Platform
Team), the researcher's dual role as participant observer and facilitator, the data
collection procedures (participant observation, artifact analysis, and open-ended active
listening sessions), the thematic coding analysis approach, and the criteria used to
establish trustworthiness and ethical conduct.

Chapter~\ref{chap:results} presents the empirical results and refinements of the study.
The chapter follows the research questions in order: how teams used Hypo-Stage and
ArchHypo in practice (RQ1); the architectural uncertainties and hypotheses they
recorded, including aggregate data characterization and Group~B structural patterns
(RQ2); the benefits and challenges they perceived (RQ3); and the refinements to
Hypo-Stage and its usage guidelines that emerged from the study, each traced to
observations and linked primarily to (RQ4). Presenting results and refinements together
makes explicit how empirical findings have driven iterative tool evolution throughout the
observation period.

Finally, Chapter~\ref{chap:conclusion} builds on the integrated empirical account in Chapter~\ref{chap:results}. 
It synthesizes findings in relation to the research questions and to prior ArchHypo evidence, discusses implications for practitioners adopting
hypothesis-engineering tools and for researchers studying tool-mediated architectural
practices, and outlines limitations together with directions for future work on Hypo-Stage and on broader
support for distributed architectural decision-making in agile organizations.

%!TeX root=../thesis.tex

\chapter{Background}
\label{cap:background}
\label{chap:background}
\label{ch:background}

This chapter establishes the theoretical foundations for the empirical investigation
carried out in this work. It begins with an overview of software architecture and
architectural uncertainty, then introduces hypotheses engineering and the ArchHypo framework in technical details, describes Internal
Developer Portals and their role as platforms for enabling tools such as Hypo-Stage,
presents the Hypo-Stage tool itself, and concludes with a summary of findings from the systematic mapping study that mapped the gaps between architectural derivation methods and practical tool support \citep{souza2019}.

% ---------------------------------------------------------------------------
\section{Software Architecture and Architectural Uncertainty}
\label{sec:sw-arch}

Software architecture refers to the high-level structure of a software system, encompassing its components, their relationships, and the design principles that guide
its evolution ~\citep{bass2021}. However, in organizations that maintain software at the core of their business developed and operated by multiple teams,
the ever-happening architectural decisions shape not only technical quality attributes such as performance, reliability, and maintainability of systems, 
but also the ability of engineering teams to work independently and deliver changes quickly and efficiently at a speed that needs to match the speed of the business ~\citep{skelton2019}.

More embracing to that rapidly changing business scenarios, \citet{richardsFord2020} added that the design decisions themselves are part of the software architecture itself, 
making it not only a static structure or frozen blueprint of a system, but a living body of reasoning that is constantly being refined and adapted to the changing business needs. 
This framing is adopted in this work to ground the concept of Architectural Uncertainty: the rationale behind any architectural choice is valid only relative to the information available at the time of the decision, 
and that information is always incomplete and subject to change.

In other words, Architectural Uncertainty arises whenever teams must make structural decisions under
incomplete information. This is a normal condition in continuous-delivery environments
where business requirements, user behavior, and technology options are only partially
known. When uncertainty is left implicit, teams either defer decisions indefinitely
(accumulating technical debt) or make premature actions that may constrain the
architecture evolution of their systems. Both outcomes can slow down delivery and degrade system
quality over time, turning the software architecture into what Foote and Yoder \citeyearpar{footeYoder1999} famously characterized as a ``Big Ball of Mud.''

Moreover, the concept of ``fast flow'', popularized through Team Topologies \citep{skelton2019}, emphasizes that organizations should design both their
teams and their architecture to enable a continuous, rapid stream of software changes.
Achieving fast flow presupposes that teams can make local architectural decisions
without waiting for a central architect to approve them. This in turn demands that architectural
uncertainties be made explicit and that lightweight mechanisms exist for teams to assess, prioritize, and revisit them as evidence accumulates. 
This is where Hypotheses Engineering may be used to play a key role in the architecture work, as proposed by \citet{silva2024tse}.

% ---------------------------------------------------------------------------
\section{Hypotheses Engineering and ArchHypo}
\label{sec:archhypo}

Hypotheses engineering is an approach to managing technical uncertainty by formulating
architectural decisions as testable hypotheses rather than settled facts. Instead of
asserting that a particular technology or pattern \emph{will} satisfy a quality
requirement, teams state that it \emph{might} do so and then design experiments to
gather evidence to validate or invalidate the hypothesis.

The ArchHypo technique, proposed by \citet{silva2024tse}, operationalized this idea for
software architecture. Grounded in Design Science Research (DSR)
and evolved from a body of patterns for agile architecture first disseminated at
conferences and through corporate training reported in~\citep{silva2024ieeesw}, ArchHypo provides
a structured vocabulary and a set of practices that fit within iterative development
workflows where architecture uncertainty is a constant challenge for rapidly-evolving software systems.

% ------- 2.2.1
\subsection{Core Concepts}
\label{subsec:archhypo-core}

An \emph{architectural hypothesis} is a falsifiable statement that captures an uncertainty relevant to the design of a software architecture. 
It is important to distinguish the concept from adjacent concepts:

\begin{itemize}
  \item \textbf{Hypothesis versus\ Risk:} A risk involves a known event with estimable
  probability and impact. An architectural hypothesis, by contrast, stems from a \emph{lack of information} sufficient to make a decision. 
  The distinguishing feature is the epistemic gap, not the probability of an outcome as observed in \citet{silva2024tse}.

  \item \textbf{Hypothesis versus\ Requirement:} A requirement specifies what the system
  must do or how well it must do it. A hypothesis may be triggered by a requirement
  (e.g., a quality-attribute goal that cannot yet be confirmed as achievable), but the
  hypothesis captures the uncertainty about the \emph{path} to satisfying that
  requirement, not the requirement itself.

  \item \textbf{Hypothesis versus\ Architectural Decision.} An architectural decision
  resolves a design choice. A hypothesis \emph{precedes and informs} decisions: it
  identifies what must be learned or validated before a decision can responsibly be
  made (~\citep{silva2024tse}).
\end{itemize}

\begin{table}[htbp]
\centering
\caption{Distinguishing architectural hypotheses from adjacent concepts \citep{silva2024tse}.}
\label{tab:hypothesis-concepts}
\begin{tabular}{p{0.22\linewidth} p{0.35\linewidth} p{0.35\linewidth}}
\toprule
\textbf{Concept} & \textbf{Definition} & \textbf{Relation to Architectural Hypothesis} \\
\midrule
Risk & A known event with estimable probability and impact. & A hypothesis stems from an \emph{epistemic gap} --- lack of information --- not from a known probability. \\
\midrule
Requirement & Specifies what the system must do or how well it must do it. & A requirement may \emph{trigger} a hypothesis when the path to satisfying it is uncertain; they are not the same artifact. \\
\midrule
Architectural Decision & Resolves a design choice. & A hypothesis \emph{precedes and informs} a decision; it captures what must be learned before a decision can responsibly be made. \\
\bottomrule
\end{tabular}
\end{table}

Hypotheses arise from two principal sources:

\begin{itemize}
  \item (a)~\emph{requirements}: when quality attribute targets are poorly specified or the team lacks information about whether a given approach can satisfy them; and 
  \item (b)~\emph{solutions}: when the consequences of an existing or candidate architectural solution are unknown.
\end{itemize}

% ------- 2.2.2
\subsection{Assessment}
\label{subsec:archhypo-assessment}

Each hypothesis is evaluated along two independent dimensions:

\begin{enumerate}
  \item \textbf{Uncertainty level.} How far the team is from being able to refute or
  accept the hypothesis. Crucially, uncertainty level is \emph{not} equivalent to
  probability. A hypothesis that is almost certainly true and a hypothesis that is almost
  certainly false both represent \emph{low} uncertainty, because in either case the team
  already has enough information to act. High uncertainty means the team genuinely
  does not know which way the evidence will fall.

  \item \textbf{Impact level.} The potential consequences --— in terms of rework,
  schedule, cost, or quality degradation --— if the hypothesis is invalidated. A high
  uncertainty combined with a high impact signals that a hypothesis demands urgent
  attention.
\end{enumerate}

Both dimensions can be assessed using a 3-point or 5-point qualitative scale;
assessments are made independently and relatively, not on absolute probability as reported in \citet{silva2024tse}.

% ------- 2.2.3
\subsection{Technical Plans}
\label{subsec:archhypo-technical-plans}

For hypotheses with significant uncertainty or impact, teams formulate a
\emph{technical plan}: a set of concrete actions whose goal is to (a)~accept or refute
the hypothesis, (b)~reduce its uncertainty level, (c)~reduce its potential impact, or
(d)~define criteria under which the decision can responsibly be
deferred. ArchHypo identifies the following plan strategies as observed in \citet{silva2024ieeesw}:

\begin{itemize}
  \item \textbf{Architectural Spike.} A time-boxed, isolated experiment to explore a
  technology or design option without coupling the result to the production codebase.
  The spike produces information rather than deliverable functionality.

  \item \textbf{Tracer Bullet.} A thin, end-to-end implementation within the actual
  system, exercising all architectural layers to validate an integration or performance
  characteristic under realistic conditions.

  \item \textbf{Software Analytics.} Using monitoring, profiling, or usage data already
  collected from the running system to gather evidence about a hypothesis without
  additional development effort.

  \item \textbf{Modularity\,/\,Flexibility.} Restructuring the affected subsystem to
  reduce coupling or increase replaceability, thereby lowering the \emph{impact} of the
  hypothesis rather than resolving its uncertainty. This strategy is useful when full
  resolution is premature but protection against a bad outcome is warranted.

  \item \textbf{Guideline Implementation.} Tailoring the development process—for
  example, adopting a coding standard, a review protocol, or a testing policy—so that
  the team systematically generates evidence about the hypothesis as a by-product of
  normal work. This was the most frequently used strategy in the evaluation reported by \citet{silva2024tse}.

  \item \textbf{Architectural Trigger.} A condition-based deferral in which the team
  agrees to act on the hypothesis only when a specific observable event occurs (e.g.,
  when the number of concurrent users exceeds a threshold, or when a performance
  regression appears in monitoring). Until the trigger fires, the hypothesis is
  deliberately deferred.

  \item \textbf{Agile Landing Zones.} Adopting architectural patterns that explicitly
  preserve flexibility, creating ``landing zones'' into which future decisions can be
  inserted without major disruption, regardless of which direction the evidence
  eventually points.
\end{itemize}

\begin{table}[htbp]
\centering
\caption{ArchHypo technical plan strategies and their primary effect \citep{silva2024ieeesw,silva2024tse}.}
\label{tab:technical-plans}
\begin{tabular}{p{0.25\linewidth} p{0.35\linewidth} p{0.30\linewidth}}
\toprule
\textbf{Strategy} & \textbf{Description} & \textbf{Primary Effect} \\
\midrule
Architectural Spike & Time-boxed isolated experiment; produces information, not deliverable functionality. & Reduces uncertainty. \\
\midrule
Tracer Bullet & Thin end-to-end implementation in the real system to validate integration or performance. & Reduces uncertainty under realistic conditions. \\
\midrule
Software Analytics & Uses existing monitoring or profiling data to gather evidence. & Reduces uncertainty without extra development effort. \\
\midrule
Modularity / Flexibility & Restructures the subsystem to reduce coupling or increase replaceability. & Reduces impact (not uncertainty). \\
\midrule
Guideline Implementation & Tailors the development process (coding standards, review protocols) to generate evidence as a by-product of normal work. & Reduces uncertainty incrementally; most frequently used \citep{silva2024tse}. \\
\midrule
Architectural Trigger & Defers action until a specific observable condition is met. & Manages timing of decision; reduces premature cost. \\
\midrule
Agile Landing Zones & Adopts patterns that preserve flexibility for future insertions. & Reduces impact proactively. \\
\bottomrule
\end{tabular}
\end{table}

% ------- 2.2.4
\subsection{Lifecycle}
\label{subsec:archhypo-lifecycle}

Hypotheses progress through a defined lifecycle: they begin in an \emph{open} state
and transition to either \emph{validated} or \emph{invalidated} as evidence from
technical plans accumulates. The \emph{responsible moment} is the key decision point in the hypothesis lifecycle:
it is the point at which the team has accumulated enough evidence to act on the
architectural decision --- neither so early that the decision would be premature, nor
so late that deferral would create unnecessary risk.

% ------- 2.2.5
\subsection{Empirical Evaluations}
\label{subsec:archhypo-evaluations}

\paragraph{IEEE Software evaluation (Catch Solve startup).}
The first evaluation of ArchHypo in a real operating environment was reported in an
IEEE Software article in \citet{silva2024ieeesw}. The setting was Catch
Solve, a startup offering web testing and monitoring services based in Italy's South
Tyrol region. Four architectural hypotheses were formulated and
tracked over approximately ten months, covering quality attributes of availability,
reusability of test templates, usability, and scalability of virtual machines. For
example, the hypothesis \emph{``The lack of redundancy can cause problems with
application availability for the customers''} was assessed with a low impact and very
low uncertainty, which led the team to defer investment in redundant infrastructure and
instead set up an Architectural Trigger based on the number of new
customers. The study showed that naming and tracking
uncertainties explicitly helped the technical lead decide when to invest in architectural
improvements versus when to wait, and validated the overall structure of the ArchHypo
approach. A key insight from this evaluation was that explicit hypothesis management
changed how the team prioritized technical work: rather than treating architectural
uncertainty as urgent and tacit, it became visible, assessed, and deliberately managed
within a defined lifecycle.

\paragraph{TSE evaluation (PropSEP at Jetsoft).}
The larger and more controlled evaluation, published in the IEEE Transactions on
Software Engineering ~\citep{silva2024tse}, examined the application of ArchHypo within
Jetsoft, a Brazilian software development company, on the PropSEP project---the budget
proposal system of the Planning and Management Secretariat of the State of São Paulo.
PropSEP is mission-critical: the approved budget it processes amounts to approximately
317~billion~BRL ~\citep{silva2024tse}. The evaluation ran across six sprints of seven
business days each (May--August 2022) and involved eight participants (P1--P8), all
holding computer science degrees with certifications in Oracle, WebSphere, or GeneXus.
Ten architectural hypotheses were identified, spanning six quality attributes
(performance, security, reliability, flexibility, usability, and team productivity), and
seven distinct types of architectural practice were employed in the resulting technical
plans~\citep{silva2024tse}.

The study produced nine key findings:

\begin{enumerate}
  \item Some hypotheses carried uncertainty that was monitored throughout the entire
  project duration without reaching full resolution, demonstrating that the technique
  accommodates long-horizon uncertainties ~\citep{silva2024tse}.

  \item ArchHypo enabled the team to avoid upfront architectural design, distributing
  structural decisions across iterations as evidence became available ~\citep{silva2024tse}.

  \item Participants reported improvements to the development process in terms of
  safety, transparency, and predictability ~\citep{silva2024tse}.

  \item Hypotheses spanning different quality attributes were handled using qualitatively
  different technical strategies, illustrating the creativity the technique
  affords~\citep{silva2024tse}.

  \item Guideline Implementation (process tailoring) emerged as the most frequently
  employed strategy across all technical plans~\citep{silva2024tse}.

  \item The team collectively recognized that managing hypotheses improved the quality
  of their architectural decision-making ~\citep{silva2024tse}.

  \item Post-release evidence pointed to a positive impact on product quality: no issues
  related to the managed hypotheses were detected after deployment ~\citep{silva2024tse}.

  \item \textbf{Changes in the team's way of working constituted a significant adoption
  obstacle.} Mixing architectural and hypothesis-related tasks with functional
  development created estimation difficulties, and participants described the additional
  work as not feeling directly connected to product delivery ~\citep{silva2024tse}.

  \item \textbf{All eight participants unanimously considered the technique hard to
  learn.} Specific difficulties included: mapping architectural concerns to the hypothesis
  format, formulating appropriate technical plan actions, articulating the assumptions
  behind a hypothesis, and understanding hypotheses as the primary driver of
  architectural action. No participant felt prepared to apply ArchHypo without
  continued expert support after the initial workshop ~\citep{silva2024tse}. 
  One participant explicitly stated that *''more training would make the team less dependent on the team of architects''* (P3) ~\citep{silva2024tse}. 
  The authors themselves acknowledge this structural dependency in their discussion. 
  The absence of a supporting tool that encodes this knowledge into the workflow risks ArchHypo to remain a technique that only functions when an experienced architect is actively steering its application.
\end{enumerate}

Findings ~8 and ~9 directly motivate the tooling contribution of this work: a tool that
embeds ArchHypo scaffolding in the daily engineering workflow, reducing the cognitive
overhead of manual application and knowledge transfer.

\begin{table}[htbp]
\centering
\caption{Summary of prior empirical evaluations of ArchHypo.}
\label{tab:archhypo-evaluations}
\begin{tabular}{p{0.18\linewidth} p{0.18\linewidth} p{0.18\linewidth} p{0.18\linewidth} p{0.18\linewidth}}
\toprule
\textbf{Study} & \textbf{Setting} & \textbf{Scope} & \textbf{Key Finding} & \textbf{Limitation} \\
\midrule
\citep{silva2024ieeesw} & Catch Solve startup (Italy) & 4 hypotheses, 10 months & Explicit hypothesis tracking changed how team prioritized architectural investments. & Single practitioner; no tool support. \\
\midrule
\citep{silva2024tse} & PropSEP / Jetsoft (Brazil) & 10 hypotheses, 6 sprints, 8 participants & ArchHypo distributed architectural work across iterations; all participants found technique hard to learn. & Manual application; expert steering required; no supporting tool. \\
\midrule
This work & Large-scale tech organization (Latin America) & Multiple teams, scaled product & To be reported --- first evaluation in a tool-mediated setting via Hypo-Stage. & Single organizational context; Backstage dependency. \\
\bottomrule
\end{tabular}
\end{table}

% ---------------------------------------------------------------------------
\section{Internal Developer Portals and Backstage}
\label{sec:idp-backstage}

The adoption of microservice architectures and the pursuit of team autonomy through
practices such as Team Topologies have created significant complexity in the daily
workflow of software engineering teams. In response, the industry has converged around the
concept of \emph{Internal Developer Portals} (IDPs) as a platform-engineering
discipline aimed at reducing this complexity while preserving team
independence ~\citep{correa2025}.

% ------- 2.3.1
\subsection{The Role of Internal Developer Portals (IDPs)}
\label{subsec:idp-role}

An Internal Developer Portal serves as the primary interface between developers and the
underlying platform infrastructure. Unlike a simple wiki or a project-management tool,
an IDP is an operational hub that aggregates tools, services, documentation, and data
into a unified ``pane of glass'' ~\citep{correa2025}. Its primary goal is to reduce
\emph{cognitive load}: by abstracting the complexity of infrastructure and standardizing
workflows, IDPs enable what practitioners call \emph{Golden Paths} -- recommended,
supported, and automated paths for building and deploying software ~\citep{correa2025}.

This goal aligns directly with the theory of autonomous stream-aligned teams summarized in Team Topologies guidance \citep{skelton2019}: if teams are to own full slices of a product
without depending on a central gatekeeper, the platform must lower the barriers to
acting on that ownership. An IDP operationalizes this by ensuring that the information
and tooling a team needs—service catalog, documentation, CI/CD pipelines, monitoring
dashboards, and domain-specific plugins—are accessible from a single environment,
without context-switching.

Conceptually, IDPs sit at the intersection of two ideas that are highly relevant to this
work. First, they embody the principle of \emph{decentralized architecture}: rather
than requiring teams to escalate architectural concerns to a central architecture group,
a well-designed IDP makes it possible to encode architectural guidance, constraints, and
tools directly into the developer workflow. Second, they are the natural delivery vehicle
for hypothesis-engineering support: because ArchHypo must become part of the daily
engineering cadence --- not a separate off-sprint activity ---embedding it in an IDP portal where
developers already spend their working time is strategically motivating.

% ------- 2.3.2
\subsection{Backstage: Open Platform for Internal Developer Portals}
\label{subsec:backstage}

Backstage\footnote{\url{https://github.com/backstage/backstage}} is an open-source framework for building Internal Developer Portals,
originally developed at Spotify to manage the complexity arising from rapid growth and
an expanding microservice estate. Spotify's internal version of Backstage serves as the
central nervous system for its infrastructure, used by over 2,700 engineers to manage
more than 14,000 software components~\citep{backstage2024}. Spotify open-sourced the
framework in March~2020, and it was subsequently donated to the Cloud Native
Computing Foundation (CNCF) as an Incubating project ~\citep{backstage2024}. As of
2025, more than 3,000 companies have adopted Backstage for their own
portals ~\citep{backstage2024}.

Backstage is not a rigid application but rather a modular ecosystem driven by
configuration. Its architecture rests on three layers: (1)~the \emph{Core Framework},
which manages essential utilities including the plugin loader, routing, authentication,
and error handling; (2)~the \emph{Plugin Ecosystem}, built around a split-plugin model
where frontend plugins handle views and UI components while backend plugins manage
API services, data processing, and storage; and (3)~\emph{External Integrations}, through
which backend plugins serve as gateways to external systems and APIs, allowing the
portal to aggregate data from cloud providers, CI/CD tools, and other
infrastructure ~\citep{correa2025}.

\begin{table}[htbp]
\centering
\caption{Backstage architecture layers and their role \citep{correa2025}.}
\label{tab:backstage-layers}
\begin{tabular}{p{0.25\linewidth} p{0.35\linewidth} p{0.32\linewidth}}
\toprule
\textbf{Layer} & \textbf{Responsibilities} & \textbf{Relevance to Hypo-Stage} \\
\midrule
Core Framework & Plugin loader, routing, authentication, error handling. & Provides the lifecycle hooks that Hypo-Stage's frontend and backend plugins register into. \\
\midrule
Plugin Ecosystem & Frontend plugins (views, UI); backend plugins (API services, data, storage). & Hypo-Stage is implemented as a split plugin: the frontend renders hypotheses and assessments; the backend persists data and exposes the API. \\
\midrule
External Integrations & Backend plugins as gateways to external systems (cloud, CI/CD, APIs). & Hypo-Stage integrates with the Backstage catalog to auto-associate hypotheses with software entities. \\
\bottomrule
\end{tabular}
\end{table}

Crucially, the plugin architecture allows organizations to extend the portal's
functionality by integrating custom tools directly into the developer's daily workflow,
without forcing context-switching. It is precisely this extensibility that makes Backstage
the appropriate delivery vehicle for the Hypo-Stage tool: a plugin can introduce a new class of
first-class object -- the architectural hypothesis -- into the same environment where a
developer already inspects service health, reviews documentation, and triggers CI/CD
runs \citep{correa2025}.

\citet{correa2025} demonstrated the technical feasibility of building such a plugin
and validated the integration approach. That work established the initial version of
the Hypo-Stage plugin for Backstage, which constitutes the artifact under empirical evaluation in this work.

% ---------------------------------------------------------------------------
\section{Hypo-Stage: The Initial Tool}
\label{sec:hypo-stage}

Hypo-Stage is an open-source Backstage plugin developed
as an undergraduate capstone project by Pedro Henrique Mariano Corrêa, supervised by
the author of this work and co-supervised by Prof.\ Paulo
Meirelles \citep{correa2025}. Its purpose is to operationalize ArchHypo within the Spotify's IDP Backstage, embedding hypothesis management into the engineering platform workflow.

The initial version of Hypo-Stage, corresponding to the code up to commit
\texttt{ceee509776508081ebdd473c2c4f710b8ef55947} in the \texttt{hypo-stage}
repository,\footnote{Source repository: \url{https://github.com/ArchHypo/hypo-stage}. The baseline revision for the initial tool design referenced throughout this thesis is commit \texttt{ceee509776508081ebdd473c2c4f710b8ef55947}.} provides the
following capabilities:

\begin{itemize}
  \item \textbf{Hypothesis management}: creating, editing, listing, and deleting
  architectural hypotheses, each linked to a Backstage catalog entity (e.g., a
  service, component, or system). Integration with the Backstage catalog means that
  hypotheses are automatically associated with the software entities they affect, without
  requiring separate registration.

  \item \textbf{Uncertainty and impact assessment}: rating each hypothesis on 5-point
  Likert scales for both uncertainty level and impact level, with free-text rationale
  fields. The use of a 5-point scale rather than a binary or 3-point scale provides
  finer-grained prioritization while remaining tractable for teams without formal risk
  management backgrounds.

  \item \textbf{Technical planning}: attaching one or more technical plan items to a
  hypothesis, specifying action type (spike, tracer bullet, trigger, etc.), description,
  target date, expected outcome, and documentation links.

  \item \textbf{Visualization}: a chart plotting hypotheses on an
  uncertainty-vs-impact plane to support prioritization, making the relative urgency of
  hypotheses visible at a glance.

  \item \textbf{Status lifecycle}: tracking each hypothesis through the defined lifecycle
  states (open~$\rightarrow$~validated or invalidated), preserving the history of status
  transitions to support retrospective analysis.
\end{itemize}

By integrating with Backstage, Hypo-Stage leverages the portal's service catalog so
that hypotheses are automatically associated with the software entities they affect. This
design reduces context-switching: developers manage architectural hypotheses in the
same environment where they consult documentation, inspect service health, and access
CI/CD pipelines.

\paragraph{Known limitations:} The initial tool design assumed familiarity with the
ArchHypo vocabulary, which poses a barrier for teams new to hypothesis engineering.
The case study \citep{correa2025} reconstructed the Catch Solve startup scenario
digitally but did not involve real-time use by a team. Furthermore, the tool lacked
features for collaborative assessment (e.g., multiple raters per hypothesis) and for
tracking the evolution of assessments over time. These gaps also motivated the empirical
evaluation and refinement effort promoted in this work.

\begin{table}[htbp]
\centering
\caption{Hypo-Stage initial capabilities and known limitations identified prior to this study \citep{correa2025}.}
\label{tab:hypostage-capabilities}
\begin{tabular}{p{0.30\linewidth} p{0.32\linewidth} p{0.30\linewidth}}
\toprule
\textbf{Capability} & \textbf{Description} & \textbf{Known Limitation} \\
\midrule
Hypothesis management & Create, edit, list, delete hypotheses linked to Backstage catalog entities. & Assumed prior ArchHypo vocabulary familiarity; no onboarding guidance. \\
\midrule
Uncertainty \& impact assessment & 5-point Likert scales per dimension with free-text rationale. & No collaborative multi-rater assessment; evolution of assessments not tracked over time. \\
\midrule
Technical planning & Attach plan items with action type, description, target date, and outcome. & Action types not enforced or explained in-tool; required expert facilitation to use correctly. \\
\midrule
Visualization & Uncertainty-vs-impact scatter plot for prioritization. & Static; not updated reactively as assessments changed. \\
\midrule
Status lifecycle & Tracks hypotheses through open $\rightarrow$ validated / invalidated states. & No history of intermediate state transitions preserved. \\
\bottomrule
\end{tabular}
\end{table}

% ---------------------------------------------------------------------------
\section{Deriving Architecture from Requirements: The Tooling Gap}
\label{sec:tooling-gap}

The systematic mapping study \citep{souza2019} examined
the existing body of work on methods for deriving architectural models from requirements
specifications. Starting from 2,575 candidate papers, the authors applied a PICOC
(Population, Intervention, Comparison, Outcome, Context) search structure and a quality
assessment checklist -- composed of general criteria (weighted at 25\%) and
method-specific criteria (weighted at 75\%) -- to select 39 high-quality primary studies
for detailed analysis. Their analysis employed the NIMSAD (Normative Information Model-based Systems Analysis and Design) framework to characterize the
structure of each reviewed method.

Among their key findings, these were directly relevant to this work:

\begin{enumerate}
  \item \textbf{Dependence on tacit knowledge.} Every reviewed method relied, to some
  degree, on the architect's experience and intuition. No method provided a fully
  automated or guided derivation path accessible to less-experienced practitioners~\citep{souza2019}.

  \item \textbf{Limited support for less-experienced practitioners.} Methods assumed a
  skilled architect capable of navigating trade-offs among quality attributes. Teams without senior architectural expertise had little actionable guidance~\citep{souza2019}.

  \item \textbf{Weak or absent tool support.} Although some tools were identified—most
  notably QuaDAI and Monarch—they addressed only fragments of the derivation process
  and offered only partial coverage of the architectural reasoning cycle. The mapping
  found no comprehensive tool integrating hypothesis formulation, assessment, and
  tracking within a team's daily workflow~\citep{souza2019}.

  \item \textbf{Insufficient evaluation.} Many proposed methods lacked rigorous
  empirical validation, relying on toy examples or author-run case studies rather than
  independent evaluations with real teams in production
  environments~\citep{souza2019}.
\end{enumerate}

The mapping study also reported that 76.9\% of the 39 primary studies addressed functional
requirements, with a significant proportion also treating non-functional requirements;
yet the methods for handling quality attributes — precisely the domain where architectural
uncertainty is most consequential — were among the least well supported by tools or
empirical evidence \citep{souza2019}.

These findings reinforce the motivation for this work on two fronts. First, they confirm
the existence of a ``tooling void'' in which architectural management remains
disconnected from the daily engineering workflow. Second, they underscore the need for
empirical evaluations with real teams -- precisely the kind of evaluation this work aims
to provide for the Hypo-Stage tool and ArchHypo.

%!TeX root=../thesis.tex

\chapter{Related Work}
\label{cap:related-work}
\label{chap:related-work}

This chapter synthesizes the literature relevant to the empirical evaluation of Hypo-Stage. 
Rather than presenting a standalone systematic review, it draws on the author's earlier exploratory work --- which included a tertiary systematic
mapping study on software architecture for startups and a multivocal literature review
on centralization and decentralization of architectural decision-making --- and
repositions those findings as narrative context for this work. The chapter also
incorporates references from the gray literature. Such sources
are explicitly identified as relevant, despite having weaker evidential
weight than peer-reviewed publications, because they often capture industry trends and real
organizational experience more rapidly than academic publication cycles allow.

% ---------------------------------------------------------------------------
\section{Centralization and Decentralization of Architectural Decisions}
\label{sec:central-decentral}

A recurring theme in both academic and gray literature is the tension between
centralizing architectural authority in a few experienced individuals versus distributing it
across development teams. Organizations pursuing fast flow and team autonomy
increasingly favor decentralized models as observed by ~\cite{salameh2021architecture, deDiana2011arch},
yet several studies report that without explicit governance mechanisms, decentralization
leads to inconsistent decisions and architectural drift as seen in the work of ~\cite{silveiraNeto2025agile, dasanayake2017collaborative, deDiana2011arch}.

The author's preliminary exploratory work examined this tension through two complementary
lenses. Firstly, the author conducted a tertiary systematic mapping study surveying the existing landscape of secondary studies on software architecture for
startup and scale-up organizations, identifying recurring themes in how early-stage
teams approach architecture governance under resource and time constraints. For the second lens, the author 
performed a multivocal literature review that synthesized both academic
publications and gray literature (practitioner blogs, conference talks, and industry
reports), covering three interrelated perspectives: deriving software architecture from
requirements, the emerging socio-technical architecture trend (which situates
architectural decisions within organizational and team structures), and the tension
between centralized and decentralized decision-making.

The multivocal review surfaced a framing where the industry advocates strongly for decentralized
decisions to enable team autonomy and fast flow \cite{skelton2019}, yet the methods and tools available
for architectural decision-making depend heavily on the tacit knowledge of experienced
architects as observed by \cite{souza2019}. Decentralization thus risks transferring decision responsibility to
practitioners who lack the conceptual vocabulary or methodological support to exercise
it effectively. However, the author recognized that this tension reflects a recognized design trade-off 
rather than a genuine logical contradiction, and framing this work around resolving that tension would risk over-extending its scope.
For this reason, this work consequently does not employ that framing; instead, it focuses on the
concrete, tractable problem of providing tool support for one well-defined uncertainty
management technique.

% ---------------------------------------------------------------------------
\section{Socio-Technical Architecture}
\label{sec:sociotechnical}

A significant body of work, both academic and practitioner-oriented, frames
architectural decision-making not merely as a technical concern but as a
\emph{socio-technical} one --- shaped by the organizational structure,
communication patterns, and cognitive constraints of the teams
involved, as observed in ~\cite{tang2017human, razavian2019empirical, herbsleb1999splitting}.
This framing is directly relevant to the motivation for the Hypo-Stage tool,
because tool-embedded scaffolding does not only lower the technical cost of
applying the ArchHypo technique, but it also contributes to changing the social
dynamics of how architectural uncertainty is surfaced and addressed within a
team ~\cite{borowa2023rationales, razavian2019empirical}.

\subsection{Conway's Law and Its Implications}
\label{subsec:conways-law}

The foundational observation underlying the socio-technical framing is Conway's
Law ~\cite{conway1968}: ``Any organization that designs a system (defined broadly) will
produce a design whose structure is a copy of the organization's communication
structure.'' Conway's Law implies that architectural decisions cannot be analyzed in
isolation from the organizational and team structure in which they are made. Considering that perspective, 
this work also supports the idea that team boundaries become architectural boundaries.

~\cite{skelton2019} extended this insight into an actionable framework
with the concept of the \emph{Inverse Conway Maneuver}: rather than allowing the
organization's structure to passively determine the architecture, teams should
deliberately design team topologies to produce the target architecture. In practical
terms, this means that architectural governance --- including how uncertainty is managed
and who is empowered to make decisions --- must be considered a design decision with
organizational as well as technical consequences.

\subsection{Socio-Technical Architecture as an Industry Trend}
\label{subsec:sociotechnical-trend}

The InfoQ Software Architecture and Design Trends Report for 2024 by ~\cite{infoq2024trends}
identified \emph{socio-technical architecture} as a distinct ``Early Adopter'' trend,
explicitly describing it as replacing the earlier framing of ``architecture as a team
sport.'' The report notes that the trend encompasses two related ideas: on one side,
the evolving role of architects as mentors and technical leaders rather than sole
decision-makers; on the other, the growing recognition --- driven in part by the
dissemination of Team Topologies~\cite{skelton2019} --- that the design of a system
must consider all the people who build, support, and maintain it ~\cite{infoq2024trends}.

It is important to note that InfoQ is a practitioner-oriented publication, not a
peer-reviewed academic venue. The Trends Report reflects the editorial judgment of an
experienced editorial board drawing on community submissions, conference talks, and
practitioner case studies. As gray literature, its findings carry weaker evidential
weight than systematic empirical studies; nonetheless, they are valuable as an indicator
of what concepts are gaining traction in industrial practice and what challenges
practitioners are actively facing, as observed in ~\cite{infoq2024trends}.

In a related InfoQ article, ~\cite{infoq2025empowering} described
a transformation journey at a software organization adopting Team Topologies and
decentralized architectural decision-making. Their account illustrates several
socio-technical dynamics that are directly relevant to this work: the initial disbelief
of teams that they are actually empowered to make architectural decisions, the flood of
decisions that follows once empowerment is accepted, the consequences teams face when
decisions have unforeseen cross-boundary impacts, and the eventual self-correction as
teams develop peer accountability mechanisms. All four phases are driven by social
factors at least as much as technical ones.

A further gray-literature source relevant to this work is the InfoQ article by ~\cite{infoq2025decentralized}, which argues that ``distributed architecture is
more social than technical.'' Tibbitts draws on Andrew Harmel-Law's concept of the
Architecture Advice Process --- in which any team can make
any architectural decision as long as it has looked for advice from the affected parties and
documented the decision in an ADR. This approach is presented as a mechanism for replacing permission-seeking
with advice-seeking. The social shift from ``who approves this?'' to ``who should I talk
to before we do this?'' is presented as a cultural unlock that enables decentralized
governance without sacrificing coherence.

% ---------------------------------------------------------------------------
\section{Tool Support for Architectural Decision-Making}
\label{sec:tool-support}

Along with the author's exploratory research, several tools and platforms have been proposed to support architectural decision-making:

\begin{itemize}
  \item \textbf{ADR tools} (e.g., adr-tools, Log4brains) focus on recording and
  versioning Architecture Decision Records (ADRs). They capture the \emph{what} and
  \emph{why} of decisions but do not model uncertainty, impact, or technical plans for
  hypotheses that are not yet decisions. ADRs document settled choices; they are not
  designed to manage the pre-decision state of uncertainty that ArchHypo addresses.

  \item \textbf{Architecture analysis tools} (e.g., ATAM-based workshops, trade-off
  analysis methods) provide structured evaluation of architectural alternatives but are
  typically heavyweight, workshop-based activities rather than continuous,
  team-embedded processes. Their value lies in one-shot or periodic analysis sessions,
  not in the ongoing, iterative management of uncertainty throughout a product's
  lifecycle.

  \item \textbf{Internal Developer Portals} such as Backstage, Port, and Cortex offer
  a catalog-centric view of services and infrastructure. While they improve
  discoverability and documentation, none of these platforms natively supports
  hypothesis engineering or architectural uncertainty management. They provide the
  integration surface that Hypo-Stage exploits, but leave the hypothesis engineering
  layer entirely absent.
\end{itemize}

Hypo-Stage differentiates itself by being explicitly designed around the ArchHypo
lifecycle: from hypothesis elicitation through assessment, technical planning, and
resolution. Its integration with Backstage places this lifecycle within the developer's
daily tooling environment. No existing tool, to the best of the author's knowledge,
provides this combination.

\begin{table}[htbp]
\centering
\caption{Comparison of tools relevant to architectural decision-making support.}
\label{tab:tools-comparison}
\begin{tabular}{p{0.20\linewidth} p{0.16\linewidth} p{0.14\linewidth} p{0.14\linewidth} p{0.14\linewidth} p{0.14\linewidth}}
\toprule
\textbf{Tool / Platform} & \textbf{Primary Purpose} & \textbf{Models Uncertainty} & \textbf{Technical Plans} & \textbf{Lifecycle Tracking} & \textbf{Team-Embedded} \\
\midrule
ADR tools (adr-tools, Log4brains) & Record settled decisions. & No & No & Partial & No \\
\midrule
ATAM / Trade-off workshops & Evaluate architectural alternatives. & Partial & No & No & No \\
\midrule
IDPs (Backstage, Port, Cortex) & Developer portal; catalog; documentation. & No & No & No & Yes \\
\midrule
Hypo-Stage (this work) & Support ArchHypo hypothesis engineering. & Yes & Yes & Yes & Yes (via Backstage) \\
\bottomrule
\end{tabular}
\end{table}

% ---------------------------------------------------------------------------
\section{Empirical Studies of Architectural Practices in Industry}
\label{sec:empirical-studies}

Industry-focused empirical studies on architectural practices remain relatively scarce
compared with tool proposals and method descriptions. Among the most relevant:

\begin{itemize}
  \item The ArchHypo evaluations by ~\cite{silva2024tse,silva2024ieeesw}
  constitute the closest antecedents to this work. They demonstrated the value of
  hypothesis engineering in two distinct settings --- a web-testing startup and a
  mission-critical government budget system --- but relied entirely on manual application
  with no tool support. This work introduces the first evaluation of ArchHypo in a
  tool-mediated setting.

  \item Studies on architectural governance in agile organizations (e.g., the work on
  SAFe and Spotify-model variants) document how companies structure roles and
  ceremonies for architecture, but rarely examine explicit uncertainty-management
  techniques.

  \item The multivocal review conducted during the author's exploratory research surfaced
  practitioner reports (blog posts, conference talks, white papers) describing ad-hoc
  approaches to architectural uncertainty. These reports corroborate the need for
  structured methods but, being gray literature, carry weaker evidential value.
\end{itemize}

Based on the findings of these related studies, this work contributes to the body of empirical evidence by evaluating a
tool-supported application of ArchHypo with multiple teams, addressing a gap in both the ArchHypo literature (which has not yet studied tool-mediated use) and the broader
architectural-practices literature (which lacks studies of hypothesis-engineering tools in scaled settings).

%!TeX root=../thesis.tex

\chapter{Research Design}
\label{ch:method}
\label{chap:research-design}

This chapter describes and justifies the research design adopted in this work.
It covers the adopted research strategy along with the conceptual framework created to fulfill the research
questions. Moreover, it covers the study context and case selection, the researcher's role, the data collection procedures, and the criteria adopted to establish the trustworthiness of the findings. 
The chapter concludes with notes on the ethical considerations that governed the study.

% ---------------------------------------------------------------------------
\section{Research Strategy}
% ---------------------------------------------------------------------------

The objective of this work was to empirically evaluate and refine Hypo-Stage,
a support tool for the ArchHypo technique, by applying it with multiple software
engineering teams in a real-world organizational setting. Answering the research
questions posed in Chapter~\ref{ch:intro} requires understanding not only \emph{what}
outcomes the tool produces, but also \emph{how} teams interact with it, \emph{why}
certain difficulties arise, and \emph{which} adjustments improve its fit with existing
practices. Following the design framework for flexible qualitative inquiry \citep{robson2016}, the purpose of this work combined with its research questions clearly pointed out the need for a \textbf{flexible research design}.

Within flexible designs, the study follows a \textbf{multiple case study strategy} as defined in methodological treatments of case-study research \citep{robson2016,yin2009}.
A case study is defined by them as an empirical investigation of a particular contemporary
phenomenon within its real-life context, using multiple sources of evidence and where
the boundary between phenomenon and context is not clearly distinct. This definition
fits the present work precisely: the phenomenon under study is the use of Hypo-Stage
by a software engineering team, and its context --- the organizational, technical, and
social environment in which the team operates --- is both relevant and inseparable
from the phenomenon itself. A purely experimental design, which would require
controlled manipulation of variables and prespecification of the design before data
collection, is not appropriate here. The number of teams is small, the settings are
naturalistic and non-manipulable, and the research purpose calls for rich,
process-oriented accounts rather than statistical inference \citep{robson2016}.

The study involves \textbf{two cases}, each consisting of a distinct software
engineering team. Using multiple cases rather than a single case follows the logic
of \emph{analytic (or theoretical) generalization} from prior methodological accounts \citep{robson2016,yin2009}: the
two cases are not selected to constitute a statistical sample of teams, but to provide
complementary and contrasting instances through which patterns of tool use,
difficulties, and refinements can be examined and cross-validated. Multiple case studies are analogous to multiple
experiments where each case either replicates findings or provides theoretically predicted contrast --- both strengthening the analytic basis for generalization \citep{robson2016}.

Beyond its descriptive and explanatory purposes, the study also has an
\textbf{evaluative purpose}: it seeks to assess the worth and fitness of Hypo-Stage
as a practical tool for managing architectural uncertainty in software engineering
teams. Flexible-design accounts distinguish \emph{formative evaluation}, which
aims at understanding and improving an ongoing program or intervention, from
\emph{summative evaluation}, which renders a final verdict on overall effectiveness \citep{robson2016}.
This study is explicitly formative: data collected during the observation period are used to identify difficulties and trigger tool refinements in near real-time, rather than to produce a retrospective pass-or-fail judgment.

Finally, the iterative loop through which the researcher observes difficulties,
implements refinements to the tool, deploys an updated version, and resumes
observation incorporates a core element of \textbf{action research} as described
in the methodological literature \citep{robson2016}: a deliberate cycle of plan, act, observe,
and reflect, with the explicit intent of producing practical change alongside
research knowledge. This does not transform the study into a full action research
project --- the primary goal remains the production of knowledge about tool use and
architectural uncertainty, not organizational change more broadly --- but the
iterative refinement component follows the same logic and must be acknowledged as
part of the design.

In summary, the research strategy can be formally described as a
\textbf{flexible, multiple case study design with formative evaluation purpose and
action research elements}, following the taxonomy from \citet{robson2016}.

\subsection{Key Methodological Concepts}
\label{sec:methodological-concepts}

Because the subsequent sections, and the results chapter, rely repeatedly on a
compact set of methodological terms drawn from flexible-design guidance \citep{robson2016}, the
definitions below clarify their meaning for the remainder of this work. They are
collected here so that every later mention resolves to an unambiguous referent.

\paragraph{Flexible design:} A research design in which the detailed plan (data
sources, instruments, and categories of analysis) is allowed to evolve during
the study in response to what the researcher learns in the field, in contrast
to a \emph{fixed design}, where all elements are pre-specified before data
collection begins~\citep{robson2016}. Flexible designs are the standard choice
for case studies, ethnographies, and grounded-theory work, where the phenomenon
under study is only partially understood at the outset and the researcher
expects to refine the analytical lens as evidence accumulates.

\paragraph{Multiple case study:} An empirical investigation that examines two
or more contemporary phenomena within their real-life contexts, using multiple
sources of evidence, where the boundary between phenomenon and context is not
clearly distinct~\citep{robson2016,yin2009}. The cases are selected not as a
statistical sample but as analytical units that either replicate findings
(\emph{literal replication}) or produce theoretically predicted contrasts
(\emph{theoretical replication}).

\paragraph{Formative evaluation:} An evaluative purpose aimed at understanding
and improving an ongoing program, intervention, or artifact while it is still
being shaped, as distinguished from \emph{summative evaluation}, which delivers
a retrospective pass-or-fail verdict on overall effectiveness~\citep{robson2016}.
In the present study, the evaluative object is the Hypo-Stage tool, and the
formative stance justifies the iterative refinement cycle described in
Section~\ref{sec:procedure}.

\paragraph{Action research:} A tradition of inquiry in which the researcher
deliberately alternates between observing a setting, intervening in it, and
reflecting on the effects of that intervention, following a recurrent
\emph{plan--act--observe--reflect} cycle~\citep{robson2016}. The present study
borrows the cycle as a design element without
adopting the emancipatory or organizational-change commitments that define
full action research.

\paragraph{Purposive sampling:} A non-probabilistic sampling strategy in which
participants or cases are chosen because they are expected to be informative
about the phenomenon of interest, rather than because they are representative
of a population~\citep{robson2016}. The two teams in this study were selected
purposively to provide complementary contrast (product vs. platform scope),
not to constitute a statistical sample of engineering teams.

\paragraph{Analytic generalization:} The form of generalization proper to case
study research, in which findings are extended to theoretical propositions
rather than to populations~\citep{robson2016,yin2009}. It is contrasted with
\emph{statistical generalization}, which requires a representative sample and
probabilistic inference. The claims in Chapter~\ref{chap:results} are analytic:
they describe conditions under which the observed patterns are theoretically
expected to recur, not the frequency with which they would appear in a wider
population.

\paragraph{Member checking:} A trustworthiness procedure in which
interpretations derived by the researcher are returned to the participants
whose behavior or statements produced them, so that the participants may
verify, correct, or elaborate on the researcher's reading~\citep{robson2016}.
Section~\ref{sec:trustworthiness} describes how member checking was
operationalized in this study through the active-listening debriefs with each
team's technical leader.

These concepts recur throughout Sections~\ref{sec:research-context},
\ref{sec:data-collection}, and~\ref{sec:data-analysis}, and frame the claims
made in Chapters~\ref{chap:results} and~\ref{chap:conclusion}.

% ---------------------------------------------------------------------------
\section{Conceptual Research Framework}
% ---------------------------------------------------------------------------

The conceptual research framework organizes the key constructs of this work and their relationships in a
way that is consistent with the research questions and guides the data collection and analysis as described for flexible designs \citep{robson2016}.
Figure~\ref{fig:conceptual-framework} describes the framework.

\begin{figure}[htbp]
  \centering
  \includegraphics[width=0.95\textwidth]{conceptual-framework}
  \caption{Conceptual framework of the study.}
  \label{fig:conceptual-framework}
\end{figure}

The central piece is the \emph{Hypo-Stage Tool Usage}, understood as the practices through which
team members appropriate and practice the ArchHypo technique via Hypo-Stage: creating
and editing architectural hypotheses, assessing uncertainty levels and impact,
defining and tracking technical plans, and revisiting hypotheses over successive
sprint cycles \citep{silva2024tse}. The tool usage is conditioned by two sets of
contextual factors: the \emph{team context} (composition, experience profile, scope of
work, and familiarity with architectural concerns) and the \emph{organizational
context} (the company's development platform, the role of Backstage as the delivery
portal, and the scaled product being developed). The tool usage produces two observable
streams: the \emph{artifacts} recorded in Hypo-Stage (hypotheses, assessments,
technical plans, and their evolution over time) and the \emph{experiential observations}
of team members (benefits, difficulties, and adaptation strategies). These streams
feed the \emph{Hypo-Stage Tool Refinements} --- changes made to Hypo-Stage and its usage
guidelines --- which in turn modify subsequent tool usage, forming an iterative cycle.
This framework is consistent with the formative evaluation and action research elements of the design, and directly links each element to the research questions (RQ1--RQ4) stated in Chapter~\ref{ch:intro}.

% ---------------------------------------------------------------------------
\section{Research Context}
\label{sec:research-context}
% ---------------------------------------------------------------------------

The study was conducted within a technology organization that operates
large-scale software products across Latin America and employs over 500 software
engineers. The organization is an intensive user of Backstage\footnote{\url{https://github.com/backstage/backstage}}, an open-source
developer portal, which serves as the central platform for internal tooling,
documentation, and engineering workflows. The availability of this platform was an
enabling condition for the study: Hypo-Stage was deployed as a Backstage plugin, so its adoption is conditioned by the availability of that platform.
The Hypo-Stage tool was deployed directly in the company's production environment, integrated with the same catalog and workflow interface that engineers use daily, reducing adoption friction considerably. 
This allowed the researcher to observe the tool usage in a naturalistic setting, with minimal interference from the researcher.

\subsection{Case Selection}
\label{sec:case-selection}

Two teams were selected as cases. Following methodological guidance on purposive sampling in flexible designs \citep{robson2016}, the cases were not chosen
to be statistically representative of a population of teams, but to provide
\emph{complementary contrast}: teams with different compositions, scopes, and
relationships to the broader system, so that differences and similarities in tool use
can be analytically compared.

\begin{description}
  \item[Team~1 --- Product Team.] A team of 12 software engineers (5 senior,
    7 mid-level) responsible for developing new product features with direct
    impact on the company's business objectives and end-user experience. This team
    works on a bounded product domain, with close proximity to business requirements
    and user feedback. Its architectural concerns are primarily about supporting new
    features reliably and at scale while preserving existing quality attributes.

  \item[Team~2 --- Platform Team.] A team of 4 software engineers (2 senior,
    2 mid-level) responsible for developing and maintaining broad-scope internal
    tooling that supports software development across the entire organization. This
    team's architectural decisions affect how all engineers in the company build,
    deploy, and operate software. Its architectural concerns are therefore more
    horizontal, with higher impact on organizational engineering practices and
    greater uncertainty about usage patterns and requirements.
\end{description}

\paragraph{Literal and theoretical replication.} In the multiple-case replication logic for case studies \citep{yin2009}, two or more cases are deliberately selected to perform one of
two complementary analytical roles. \emph{Literal replication} chooses cases
that share the relevant conditions and are therefore predicted to produce
similar findings; convergence across literal replications strengthens the
claim that the finding is genuinely tied to the intervention rather than to
idiosyncratic case features. \emph{Theoretical replication} chooses cases
that differ on a theoretically meaningful dimension and are therefore
predicted to produce contrasting findings; divergence along the predicted
direction strengthens the theoretical account of \emph{why} the intervention
behaves differently across conditions.

The two cases represent a \emph{literal replication} in so far as both teams face
the same intervention (use of Hypo-Stage and ArchHypo) and a \emph{theoretical
replication} in that their different scopes and compositions should, theoretically,
lead to different experiences of architectural uncertainty and different challenges
in adopting a hypothesis-driven approach \citep{robson2016,yin2009}. This contrast
strengthens the analytic basis for identifying which findings are context-independent and which are conditioned by team type.

\begin{table}[htbp]
\centering
\caption{Case comparison: Team 1 (Product) and Team 2 (Platform).}
\label{tab:case-comparison}
\begin{tabular}{p{0.30\linewidth} p{0.32\linewidth} p{0.32\linewidth}}
\toprule
\textbf{Dimension} & \textbf{Team 1 --- Product Team} & \textbf{Team 2 --- Platform Team} \\
\midrule
Size & 12 engineers (5 senior, 7 mid-level). & 4 engineers (2 senior, 2 mid-level). \\
\midrule
Scope & Bounded product domain; new features with direct business and user impact. & Broad-scope internal tooling supporting all engineers in the organization. \\
\midrule
Architectural concerns & Supporting new features reliably at scale; preserving existing quality attributes. & Horizontal decisions affecting build, deploy, and operate practices organization-wide. \\
\midrule
Uncertainty profile (expected) & Moderate uncertainty, high proximity to user feedback. & Higher uncertainty about usage patterns; high cross-team impact of decisions. \\
\midrule
Replication logic & Literal replication --- same intervention (Hypo-Stage + ArchHypo). & Theoretical replication --- different scope and composition expected to surface different challenges. \\
\bottomrule
\end{tabular}
\end{table}

% ---------------------------------------------------------------------------
\section{Researcher's Role}
\label{sec:researcher-role}
% ---------------------------------------------------------------------------

Throughout the study, the researcher occupied the dual role of
\textbf{participant observer and facilitator}, in the terms used in flexible-design methodology \citep{robson2016}.
The researcher is an employee of the organization in which the study was conducted, which makes this an instance of \emph{insider research}.
This position provided privileged access to the teams, their work context, and their daily practices, which would have been substantially harder to
obtain as an external investigator. It also allowed the researcher to introduce
Hypo-Stage into the company's infrastructure in a credible and practically feasible
way.

As discussed in accounts of insider research \citep{robson2016}, such studies raise specific concerns
regarding bias and the potential conflation of the researcher's own perspective with
those of participants. These concerns are addressed in
Section~\ref{sec:trustworthiness}. It is important to note that, while the researcher
served as facilitator during training sessions and provided assistance on request
throughout the observation period, he did not direct the teams in \emph{which}
architectural uncertainties to address or \emph{how} to formulate specific
hypotheses. The content of the work produced by the teams in Hypo-Stage reflects
their own judgments and priorities.

% ---------------------------------------------------------------------------
\section{Procedure}
\label{sec:procedure}
% ---------------------------------------------------------------------------

\subsection{Training and Introduction}

Before the observation period began, each team participated in a one-hour training
session led by the researcher. The session covered:

\begin{itemize}
  \item[(i)] The conceptual foundations of hypotheses engineering and the ArchHypo technique, including the distinction between architectural hypotheses, uncertainty levels, impact assessment, and technical plans.
  \item[(ii)] The context and motivation of the research, including the broader challenge of making architectural uncertainties explicit in teams working on scaled products.
  \item[(iii)] A live demonstration of Hypo-Stage running in the company's Backstage instance, illustrating how to create and manage hypotheses within the tool. Teams were given written communication prior to the meeting (Appendix~\ref{apx:training-communication}) --- including an illustration of the ArchHypo big picture (Figure~\ref{fig:archhypo-overview}) --- so that they could arrive at the session with some prior exposure to the key concepts.
\end{itemize}

\begin{figure}[htbp]
  \centering
  \includegraphics[width=0.85\textwidth]{archhypo-big-picture}
  \caption{Overview of the ArchHypo technique showing the iterative cycle through
    which engineers identify architectural hypotheses, assess their uncertainty
    and impact, define technical plans, and take or postpone architectural
    decisions. Source: adapted from \citet{silva2024tse}.}
  \label{fig:archhypo-overview}
\end{figure}

\subsection{Observation Period}

Following the training session, each team entered an observation period during which they were invited to use Hypo-Stage to
register and manage any architectural uncertainties that arose in the course of their
regular development work. The organization operates on two-week agile sprints. The initial participant communication (Appendix~\ref{apx:training-communication}) described a four-week horizon; in practice, teams continued to use the tool voluntarily beyond that window.

The empirical window for hypothesis creation is anchored in Hypo-Stage registration timestamps. As documented in Appendix~\ref{app:hypotheses-data}, the thirteen hypotheses were created between 11~February~2026 and 12~April~2026: an interval of 61~calendar days (approximately nine weeks), spanning five consecutive two-week sprint periods from the earliest to the latest registration. No constraints were imposed on the type, number, or
formulation of hypotheses, nor on when to create them during the sprints. It is intrisic to the nature of the organization in which the experience was conducted that the teams are free to adjust their agile practices to their own needs. Therefore, in contrast to the PropSEP evaluation setting \citep{silva2024tse}, the teams were not required to create hypotheses during the sprints, but rather to use the tools whenever they deemed necessary.
The objective was to observe tool use in naturalistic conditions, with minimal researcher interference in the content of the teams' work.

During this period, the researcher remained available to both teams for support on
request. Support was of two distinct types:

\begin{enumerate}
  \item \textbf{Conceptual support.} When a team member found it difficult to
    formulate a hypothesis or to assess its uncertainty and impact within the
    ArchHypo framework, the researcher provided clarification or facilitated a brief
    discussion. Interactions of this type were noted as observations about the
    conceptual difficulty of the technique.

  \item \textbf{Tool support.} When a team member encountered a difficulty that
    originated from the tool itself --- a usability problem, a missing feature, or a
    mismatch between the tool's behavior and the expected workflow --- the researcher
    logged this as a tool improvement candidate and, as a matter of priority,
    implemented and deployed an updated version of Hypo-Stage as quickly as
    practicable. This was done to ensure that tool limitations would not persistently
    obstruct teams' ability to practice ArchHypo, and to allow subsequent
    observations to reflect a more refined version of the tool.
\end{enumerate}

This second type of support operationalizes the \emph{action research cycle} within
the study: difficulty observed $\rightarrow$ refinement designed $\rightarrow$
updated tool deployed $\rightarrow$ use and observation resumed. Each refinement
cycle is documented in Section~\ref{sec:refinements}.

% ---------------------------------------------------------------------------
\section{Data Collection}
\label{sec:data-collection}
% ---------------------------------------------------------------------------

\begin{table}[htbp]
\centering
\caption{Data collection instruments, sources, and research questions addressed.}
\label{tab:data-collection}
\begin{tabular}{p{0.25\linewidth} p{0.28\linewidth} p{0.20\linewidth} p{0.19\linewidth}}
\toprule
\textbf{Instrument} & \textbf{Data Produced} & \textbf{Collection Timing} & \textbf{RQs Addressed} \\
\midrule
Participant observation (field notes) & Reactions, use patterns, informal commentary, support requests. & Continuous during observation period. & RQ1, RQ3 \\
\midrule
Artifact analysis (Hypo-Stage records) & Hypotheses, assessments, technical plans, lifecycle transitions. & End of observation period. & RQ1, RQ2 \\
\midrule
Open-ended active listening sessions & Perceived usefulness, difficulties, suggestions, overall assessment. & End of observation period. & RQ3, RQ4 \\
\bottomrule
\end{tabular}
\end{table}

Following the multiple data source approach recommended for case studies \citep{robson2016,yin2009}, data were collected through three complementary channels:

\subsection{Participant Observation}
\label{subsec:participant-observation}

The researcher's continuous presence in the organization and involvement in both
teams' environments provided an ongoing stream of observational data. Observations
covered: 

\begin{itemize}
  \item[(i)] Team members' reactions and difficulties during and after the training
  session;
  \item[(ii)] Patterns of use (or non-use) of Hypo-Stage during the sprint;
  \item[(iii)] Interactions in which team members requested conceptual or tool support;
  \item[(iv)] Informal commentary about the fit of ArchHypo with existing practices.
\end{itemize}

Following established guidance on participant observation in flexible designs \citep{robson2016}, field notes were recorded as close as possible in time to the observed events, and the researcher maintained a reflexive awareness of his dual role as insider and investigator.

\subsection{Artifact Analysis}
\label{subsec:artifact-analysis}

The hypotheses, assessments, and technical plans created by each team in Hypo-Stage
constitute a documentary record of how they used the tool over the observation period.
These artifacts were collected at the end of the observation period and analyzed to
characterize:

\begin{itemize}
  \item[(i)] The number and types of hypotheses created;
  \item[(ii)] The quality attributes and sources of uncertainty they addressed;
  \item[(iii)] How hypotheses were assessed and what technical plan actions were defined;
  \item[(iv)] How hypotheses evolved (created, refined, or closed) across the observation period (five two-week sprints spanned by registration timestamps; Appendix~\ref{app:hypotheses-data}).
\end{itemize}

The content analysis of documentary record is a recognized data collection method in flexible case study designs \citep{robson2016}.

\subsection{Open-Ended Active Listening Sessions}
\label{subsec:active-listening}

In addition to informal interactions arising from support requests, the researcher
conducted brief open-ended active listening sessions with team members at the end of the
observation period, inviting them to reflect on their experience of using Hypo-Stage
and ArchHypo. These conversations followed a semi-structured format, covering the
perceived usefulness of the technique, the difficulties encountered, the changes they
would suggest to the tool or its integration into their workflow, and their overall
assessment of the value of making architectural uncertainties explicit in this way.
Semi-structured or focused interviews are a central method for flexible design
research and are particularly appropriate when the researcher seeks to capture
participants' perspectives and interpretations rather than confirming predetermined
hypotheses as described for flexible qualitative designs \citep{robson2016}.

% ---------------------------------------------------------------------------
\section{Data Analysis}
\label{sec:data-analysis}
% ---------------------------------------------------------------------------

Data from the three collection sources described in
Section~\ref{sec:data-collection} were analyzed using \emph{thematic coding
analysis} as specified for qualitative data gathered under a flexible design \citep{robson2016}. Thematic coding is a family of qualitative analytic
techniques that produces a structured account of what is present textual or observational data by iteratively assigning short labels ---
\emph{codes} --- to meaningful segments, aggregating those codes into
higher-order \emph{categories} or \emph{themes}, and then comparing the
resulting structure across units of analysis~\citep{robson2016}. It is
distinguished from fully inductive \emph{grounded theory} by its willingness
to accept a priori analytical constructs (in this study, those of the
conceptual framework in Figure~\ref{fig:conceptual-framework}) and from fully
deductive \emph{content analysis} by its commitment to letting the codebook
emerge from the data rather than being fixed in advance.

In line with established guidance on flexible designs \citep{robson2016}, the analysis proceeded
through three stages, each answering a distinct analytical question.

\paragraph{Stage 1 --- Inductive open coding:} All field notes, extracted
Hypo-Stage artifact records, and active-listening transcripts were reviewed
line-by-line. A short descriptive \emph{code} was assigned to each segment that carried
analytical weight, capturing a recurring observation, a reported experience,
or a notable silence. Codes at this stage are descriptive, not interpretive:
the aim is faithful, low-inference labeling of what is present in the data.
No a priori categories were imposed. Representative codes and their
traceability to later sections appear in Table~\ref{tab:code-traceability}.

\paragraph{Stage 2 --- Axial grouping into conceptual categories:} The
descriptive codes produced in Stage~1 were grouped into higher-order
\emph{axial categories} that correspond to the constructs of the conceptual
framework. The term \emph{axial} follows the Strauss--Corbin tradition adopted
in accounts of thematic coding under flexible designs \citep{robson2016}: it denotes the relational re-organization of open codes
around a central analytical axis, in this study the constructs of tool usage,
architectural uncertainty, perceived benefits, perceived challenges, candidate
refinements, and organizational-adoption conditions. Six axial categories
emerged, each corresponding to one construct in
Figure~\ref{fig:conceptual-framework} or to an axis along which the
framework's observables were stratified.

\paragraph{Stage 3 --- Cross-case comparison:} The axial categories were then
examined across the two cases (Product Team and Platform Team) to identify
findings common to both teams and findings that were specific to one.
Similarities are reported as shared observations in
Chapter~\ref{chap:results}; differences are attributed to the distinct scopes
and compositions of product versus platform work, following the logic of
literal and theoretical replication set out in
Section~\ref{sec:case-selection}.

The complete mapping from representative open codes to axial categories,
research questions, and chapter sections appears in
Table~\ref{tab:code-traceability} at the start of
Chapter~\ref{chap:results}. That table is the primary audit trail linking
raw observations to every empirical claim made in this thesis.

Consistent with the observation that flexible designs may
include modest quantitative summaries without becoming multi-strategy
designs \citep{robson2016}, the artifact data extracted from Hypo-Stage were also summarized
descriptively (counts of hypotheses per team, frequency of quality
attributes, technical-plan strategy distribution). These summaries
complement and illustrate the qualitative findings; they do not constitute
a separate quantitative strand of analysis.

% ---------------------------------------------------------------------------
\section{Establishing Trustworthiness}
\label{sec:trustworthiness}
% ---------------------------------------------------------------------------

\begin{table}[htbp]
\centering
\caption{Trustworthiness criteria and strategies applied in this study \citep{robson2016}.}
\label{tab:trustworthiness}
\begin{tabular}{p{0.20\linewidth} p{0.35\linewidth} p{0.37\linewidth}}
\toprule
\textbf{Criterion} & \textbf{Meaning} & \textbf{Strategy Applied} \\
\midrule
Credibility & Internal validity; do findings accurately represent participants' realities? & Prolonged engagement ($\approx$9 weeks, five two-week sprints covered by registrations); triangulation across three data sources; member checking. \\
\midrule
Transferability & Degree to which findings may inform other settings. & Thick description of context; analytic (not statistical) generalization claimed. \\
\midrule
Dependability & Reliability and consistency of the research process. & Audit trail: dated field notes, systematic artifact extraction, commits linked to each refinement. \\
\midrule
Confirmability & Degree to which findings reflect data rather than researcher bias. & Explicit role separation; active search for disconfirming evidence; all claims grounded in specific observations or artifacts. \\
\bottomrule
\end{tabular}
\end{table}

The trustworthiness of findings from flexible design research is discussed in terms of \emph{credibility}, \emph{transferability},
\emph{dependability}, and \emph{confirmability} \citep{robson2016}. Each is addressed
below.

\textbf{Credibility} --- the internal validity of the findings --- was supported
through: 
\begin{itemize}
  \item[(i)] Prolonged engagement with both teams over approximately nine calendar weeks of real work in their natural environment (Hypo-Stage registrations span five consecutive two-week sprints; Appendix~\ref{app:hypotheses-data});
  \item[(ii)] Multiple data sources (observation, artifact analysis, active listening sessions), providing triangulation of evidence;
  \item[(iii)] \textbf{Member checking.} Interpretations that the researcher
    formed during the observation period --- for instance, the inference that
    hypothesis registration clustered around protected discussion time ---
    were surfaced back to each team's senior technical leader during the
    active-listening sessions described in
    Section~\ref{subsec:active-listening}. Each leader was explicitly invited
    to confirm, correct, or elaborate on the researcher's reading. This
    operationalizes the member-checking procedure recommended as a core credibility control for flexible qualitative
    designs \citep{robson2016}.
\end{itemize}

\textbf{Transferability} --- the degree to which the findings may inform other
settings --- is constrained by the specifics of the organizational context (a large
technology company in Latin America with an established Backstage platform) and the
tool (Hypo-Stage). The study does not claim statistical generalization. Rather, it
offers analytic generalization \citep{robson2016,yin2009}: the patterns of use,
difficulties, and refinements documented here provide a theoretically informed
account that researchers and practitioners can assess for relevance to their own
contexts.

\textbf{Dependability} --- the reliability of the research process --- was addressed
by maintaining an audit trail: field notes were dated and linked to specific events,
artifact data were extracted and stored systematically, and each tool refinement was
documented with its triggering observation and the corresponding commit in the
Hypo-Stage repository.

\textbf{Confirmability} --- the degree to which findings reflect the data rather
than researcher bias --- was addressed by: 

\begin{itemize}
  \item[(i)] Explicitly distinguishing between the researcher's facilitation role (which involved helping teams understand the technique) and his investigator role (which required non-directive observation of how they used it);
  \item[(ii)] Seeking disconfirming evidence, that is, noting instances where teams found value in the tool or the technique as readily as instances of difficulty;
  \item[(iii)] Grounding all claims in specific observations, artifacts, or conversation excerpts, reported in Chapter~\ref{chap:results} (including Section~\ref{sec:refinements} where refinements are documented).
\end{itemize}

% ---------------------------------------------------------------------------
\section{Ethical Considerations}
\label{sec:ethical-considerations}
% ---------------------------------------------------------------------------

The study raises ethical considerations associated with insider research \citep{robson2016}. Participation was voluntary: team members were informed of the research objectives, were free to use Hypo-Stage as much or as little as they chose, 
and were invited but not required to engage in active listening sessions. No personal performance data were collected or reported; all analysis
pertains to team-level patterns and tool-level observations. The organization in which the study was conducted and the teams involved are anonymized in reporting.
As no sensitive personal data were collected, no ethical review was required for this study.

The dual role of the researcher as organizational employee and academic investigator is a potential source of role conflict. This was managed by maintaining a clear
separation between the researcher's responsibilities to the organization (deploying and supporting a useful tool) and his responsibilities as a researcher (reporting
findings honestly, including negative ones). The research was conducted with the knowledge and informal consent of the relevant engineering leadership in the organization.

% ============================================================
%  Chapter 5 — Results and Refinements
%  Master's Thesis: Evaluating and Refining Hypo-Stage
%  José Gonçalves Lima Neto — IME-USP, 2026
% ============================================================
\chapter{Results and Refinements}
\label{chap:results}
\label{ch:results}

This chapter presents the empirical findings of this study and the refinements applied to
Hypo-Stage that emerged from them. The results are organized around the four
research questions defined in Chapter~\ref{chap:introduction}: 
\begin{itemize}
  \item RQ1 addresses how teams used the tool and the technique in practice;
  \item RQ2 characterizes the architectural uncertainties they registered;
  \item RQ3 reports the benefits and challenges perceived by participants;
  \item RQ4 documents each refinement, tracing it to a concrete observation and to the question it addresses.
\end{itemize}

The section~\ref{sec:emerging-patterns} goes beyond a direct answer to those four
questions by drawing two families of \textit{emerging patterns}: team-behavior
patterns (Group~A) and a set of novel technical architectural patterns for
hypothesis engineers (Group~B). Group~B constitutes a distinct contribution
of this work, addressing a gap left open by prior research in the ArchHypo line,
as discussed in Section~\ref{sec:group-b-patterns}.

\section*{Research Methodology Traceability}
\addcontentsline{toc}{section}{Methodological Traceability}

The evidence presented in this chapter was collected and analyzed strictly in
accordance with the research design described in Chapter~\ref{chap:research-design}.
This section explicitly describes this correspondence, so that we can assess
compliance with the research methodology at a glance.

Regarding data collection, three data sources were employed, as specified in Section~\ref{sec:data-collection}:
\begin{enumerate}
  \item \textbf{Participant observation:} Evidence was gathered through the
    researcher's continuous presence within the organization. Observations
    were the author's own first-hand field notes, recorded as close as
    possible in time to the observed events (Section~\ref{subsec:participant-observation}). They covered
    team reactions during and after training, patterns of tool use and
    non-use during the observation sprints, informal commentary, and all
    instances in which team members requested conceptual or tool support.
  \item \textbf{Artifact analysis:} At the end of the observation period, the
    complete set of hypotheses, assessments, and technical plans records created by each team in Hypo-Stage was extracted and analyzed
    (Section~\ref{subsec:artifact-analysis}). The quantitative summaries in Sections~\ref{sec:rq2-results} and~\ref{subsec:refinement-summary} 
    are derived exclusively from those artifact records.
  \item \textbf{Open-ended active listening sessions:} The researcher
    conducted brief, semi-structured active listening conversations with the
    senior technical leaders of each team at the end of the observation period
    (Section~\ref{subsec:active-listening}). The sessions were open-ended: each leader was invited to
    reflect on their experience of using Hypo-Stage and ArchHypo. Questions
    used in these sessions are reproduced in Appendix~\ref{apx:instruments}
    (Section~\ref{sec:listening-guide}). Literal statements from these
    sessions are reported in this chapter only when they add evidence beyond
    what is already documented in the recorded artifacts; in every case the
    participant's role is identified (e.g., ``Product Team technical leader'').
\end{enumerate}

Regarding data analysis, the three-stage thematic coding procedure defined in
Section~\ref{sec:data-analysis} was applied to the field
notes, artifact summaries, and active-listening transcripts. The stages are
summarized below for local reference; the methodological grounding is in
Section~\ref{sec:data-analysis}.
\begin{enumerate}
  \item \textbf{Stage 1 — Inductive open coding:} All primary data were
    reviewed and a set of descriptive codes was generated inductively,
    capturing recurring observations and literal themes without imposing
    prior categories. Table~\ref{tab:code-traceability} lists representative
    codes and traces each one to the axial category, research question, and
    chapter section where the corresponding finding is developed.
  \item \textbf{Stage 2 — Axial grouping into conceptual categories:} Codes
    were grouped into higher-level categories corresponding to the constructs
    of the conceptual framework (Figure~\ref{fig:conceptual-framework}).
    Six categories emerged: \emph{patterns of tool use}, \emph{types of
    architectural uncertainty and structural archetypes}, \emph{perceived
    benefits}, \emph{perceived challenges}, \emph{candidate tool refinements},
    and \emph{organizational conditions for adoption}. The right-hand columns
    of Table~\ref{tab:code-traceability} show how each category maps to a
    research question (or emerging-pattern group) and to the section of this
    chapter where the results are presented.
  \item \textbf{Stage 3 — Cross-case comparison:} Patterns within and across
    the two cases (Product Team and Platform Team) were compared to identify
    findings common to both teams and findings specific to each team's context.
    Similarities are reported as shared observations; differences are attributed
    to the distinct natures of product versus platform work, as elaborated in
    the individual sections.
\end{enumerate}

\begin{table}[htbp]
\caption{From open codes to chapter sections: traceability of the thematic coding process.}
\label{tab:code-traceability}
\centering\small
\begin{tabular}{p{0.26\linewidth} p{0.26\linewidth} l p{0.17\linewidth}}
\toprule
\textbf{Open code} & \textbf{Axial category} & \textbf{RQ / Group} & \textbf{Discussed in Section} \\
\midrule
\textit{kafka-cluster-burst-registration}
  & Patterns of tool use; Structural archetypes
  & RQ1; RQ2
  & \ref{sec:results-rq1}, \ref{sec:group-b-patterns}~(B1) \\
\midrule
\textit{senior-author-only}
  & Organizational adoption conditions
  & Group~A
  & \ref{sec:group-a-patterns}~(A1) \\
\midrule
\textit{registration-barrier-time-pressure}
  & Perceived challenges
  & RQ3
  & \ref{subsec:challenges} \\
\midrule
\textit{tool-positive-catalog-integration}
  & Perceived benefits
  & RQ3
  & \ref{subsec:benefits} \\
\midrule
\textit{falsifiability-confusion}
  & Perceived challenges; Candidate refinements
  & RQ3; RQ4
  & \ref{subsec:challenges}, \ref{sec:refinements} \\
\bottomrule
\end{tabular}
\end{table}

The artifact data from Hypo-Stage were summarized descriptively (hypothesis
counts, quality-attribute frequency, technical plan distribution) to complement
and illustrate the qualitative findings. Following methodological guidance for flexible qualitative designs \citep{robson2016}, the
inclusion of these descriptive statistics is consistent with a flexible
qualitative design and does not constitute a multistrategy approach.

% ─────────────────────────────────────────────────────────────
\section{RQ1 — How Teams Used Hypo-Stage in Practice}
\label{sec:rq1-results}
\label{sec:results-rq1}
% ─────────────────────────────────────────────────────────────

\subsection{Hypothesis Registration Dynamics}
\label{subsec:registration-dynamics}

Over the observation period (Hypo-Stage registrations span 11~February~2026--12~April~2026, approximately nine calendar weeks or five two-week sprints), both teams together registered thirteen
architectural hypotheses in Hypo-Stage. Six hypotheses originated from the
Product~Team and seven from the Platform~Team. This distribution is noteworthy
because the Platform~Team is significantly smaller (four engineers) yet produced
more registered hypotheses than the twelve-engineer Product~Team. This asymmetry
is partly explained by the registration of a dense cluster of five Platform
hypotheses on a single day (12 April 2026), all addressing complementary aspects
of the shared Kafka streaming infrastructure — a burst pattern discussed further
in Section~\ref{sec:group-b-patterns}.

Across both teams, twelve of the thirteen hypotheses were authored by senior
technical leaders — senior or staff engineers who held explicit architectural
responsibilities for their respective domains. One hypothesis was registered by a
mid-level engineer on the Product~Team. Both technical leaders in each team
independently reported that a much larger number of architectural uncertainties
existed in their daily work than the number they managed to register. The
principal constraint they named was time: formulating a well-structured hypothesis
statement requires focused, uninterrupted attention that sprint delivery
pressure rarely afforded. This observation resonates with findings from prior
ArchHypo evaluations, where the learning curve and process-integration demands
were described as significant barriers to adoption
\citep{silva2024tse, silva2024ieeesw}.

\subsection{Role-Based Engagement}
\label{subsec:role-engagement}

Participant observation and the active listening sessions revealed a clear
stratification in how different roles engaged with Hypo-Stage during the study
period. Senior engineers treated the tool as a decision-support instrument:
they used it to externalize concerns that had until then existed only as tacit
knowledge, and they referenced the hypothesis backlog when discussing
architectural trade-offs in planning sessions. Mid-level and junior engineers
were aware of the tool's existence and attended the training session, but
neither team saw them register hypotheses independently beyond the initial
introduction. This does not indicate disengagement; rather, both technical
leaders explained that these engineers typically operate at a scope where
architectural uncertainties are resolved quickly — within a sprint or less —
leaving little perceived need for formal registration.

\subsection{Tool Adoption Patterns}
\label{subsec:adoption-patterns}

Usage of Hypo-Stage was non-uniform across the observation window. Registration
activity was concentrated in two periods: a first wave during and immediately
after the training session, and a second, more intensive wave near the end of
the period when the Platform~Team used the tool for a structured architectural
discussion on Kafka failure-handling strategies. The author suggests that this
second wave is plausibly related not only to the substantive Kafka design work
but also to intensified commitment on the part of the Platform~Team's senior
technical leader to using Hypo-Stage as the agreed observation window drew toward
its close---a pattern
consistent with registrations becoming more salient when the field study's end
was imminent. This remains an interpretive inference from timing and observed
behavior, not a claim about the leader's private motives. Between these waves, both
technical leaders reported using the hypothesis detail view to review existing
records and update technical plans, but the tool was not used in team-visible
ceremonies such as sprint retrospectives or planning meetings. When asked,
participants noted that embedding the tool into these ceremonies would require
an explicit cultural agreement within the team — something that had not yet
occurred within the observation window.

The integration of Hypo-Stage into the Backstage Internal Developer Portal was
cited by both teams as a meaningful enabler of adoption. Because Backstage is
already the daily entry point for service catalog, documentation, and CI/CD
status, engineers encountered Hypo-Stage without navigating to a separate tool.
The EntityHypothesesTab, shown in Figure~\ref{fig:entity-hypotheses-tab}, which surfaces only the hypotheses linked to the
component an engineer is already viewing, was specifically mentioned as reducing
the cognitive cost of context-switching into hypothesis management.

\begin{figure}[ht]
  \centering
  \includegraphics[width=0.85\textwidth]{entity-hypotheses-tab}
  \caption{EntityHypothesesTab: Hypo-Stage surfaces only the hypotheses linked
    to the Backstage catalog component the engineer is currently viewing,
    eliminating the need to navigate to a separate tool or filter by component.}
  \label{fig:entity-hypotheses-tab}
\end{figure}

% ─────────────────────────────────────────────────────────────
\section{RQ2 — Architectural Uncertainties Registered}
\label{sec:rq2-results}
\label{sec:results-rq2}
% ─────────────────────────────────────────────────────────────

\subsection{Data Overview}
\label{subsec:data-overview}

Table~\ref{tab:hypotheses-overview} presents the thirteen hypotheses registered
during the study, anonymized in accordance with the ethical agreement described
in Section~\ref{sec:ethical-considerations}. Full exported records (statements,
ratings, entity links, and author notes) appear in Appendix~\ref{app:hypotheses-data}
(Section~\ref{sec:app-hypothesis-records}, \textit{Hypothesis Records}). Each hypothesis is identified by a
short label, its source type (Requirements or Solution), its lifecycle status
at the end of the observation period, and its uncertainty and impact ratings.

\begin{table}[htbp]
\caption{Overview of the thirteen registered architectural hypotheses.
  Req.\ = Requirements-sourced; Sol.\ = Solution-sourced.
  U = Uncertainty level; I = Impact level.
  Ratings: VL = Very Low, M = Medium, H = High, VH = Very High.}
\label{tab:hypotheses-overview}
\centering\small
\begin{tabular}{llllll}
\toprule
\textbf{ID} & \textbf{Short label} & \textbf{Team} & \textbf{Source} &
\textbf{U} & \textbf{I} \\
\midrule
H01 & Kotlin/JOOQ/Spring template readiness & Product  & Req. & VH & VH \\
H02 & Kafka direct exposure via platform framework & Platform & Req. & VL & M  \\
H03 & Decouple eng.\ logs from security/compliance audit & Platform & Req. & M  & VH \\
H04 & Blob storage GDPR right-to-erasure & Product  & Req. & M  & VL \\
H05 & Credit-bureau context-agnostic restructuring & Product  & Sol. & H  & VH \\
H06 & Billing aggregation chunked keyset pagination & Product  & Sol. & VL & VH \\
H07 & Blob ownership transfer across service domains & Product  & Req. & VL & VH \\
H08 & Blob cross-client sharing governance & Product  & Req. & VL & VH \\
H09 & DLT for audit vs.\ offset-based replay & Platform & Sol. & M  & VH \\
H10 & Synchronous retry for transient consumer failures & Platform & Sol. & M  & H  \\
H11 & Failure state inside vs.\ outside Kafka flow & Platform & Sol. & M  & H  \\
H12 & Error classification before retry & Platform & Sol. & M  & H  \\
H13 & Payload immutability before re-enqueueing & Platform & Sol. & VL & VH \\
\bottomrule
\end{tabular}
\end{table}

\subsection{Source Types}
\label{subsec:source-types}

Eight of the thirteen hypotheses were classified as \textit{Requirements}-sourced:
the uncertainty arose from an incompletely specified quality-attribute goal or
an external obligation whose implementation path is unknown. Five were classified
as \textit{Solution}-sourced: the team has an existing or candidate architectural
decision and is uncertain about its consequences. The predominance of
requirements-sourced hypotheses on the Product~Team (five out of six) and
the predominance of solution-sourced hypotheses on the Platform~Team's Kafka
cluster (four out of five Kafka hypotheses) reflects the different nature of
each team's work: the Product~Team faces open questions about whether architectural
goals can be met at all, while the Platform~Team is assessing the trade-offs of
specific implementation choices for a shared infrastructure component.

\subsection{Quality Attribute Distribution}
\label{subsec:qa-distribution}

Table~\ref{tab:qa-distribution} summarizes the frequency of quality attributes
across the thirteen hypotheses. Attributes were assigned by the hypothesis authors
using the multi-select form in Hypo-Stage; a hypothesis may carry more than one
attribute.

\begin{table}[htbp]
\caption{Quality attribute frequency across the thirteen hypotheses.}
\label{tab:qa-distribution}
\centering\small
\begin{tabular}{lcc}
\toprule
\textbf{Quality Attribute} & \textbf{Count} & \textbf{\% of hypotheses} \\
\midrule
Reliability        & 6 & 46\% \\
Auditability       & 5 & 38\% \\
Performance        & 4 & 31\% \\
Interoperability   & 4 & 31\% \\
Extensibility      & 3 & 23\% \\
Usability          & 3 & 23\% \\
Scalability        & 2 & 15\% \\
Modifiability      & 2 & 15\% \\
Flexibility        & 2 & 15\% \\
Configurability    & 2 & 15\% \\
Maintainability    & 1 &  8\% \\
Security           & 1 &  8\% \\
Observability      & 1 &  8\% \\
Reusability        & 1 &  8\% \\
\bottomrule
\end{tabular}
\end{table}

Reliability and Auditability emerge as the two dominant concerns. The prominence
of Auditability is especially notable and directly traceable to two contextual
forces: the organization's internal security audit requirements, which drove H03,
and the applicable privacy regulation governing document storage (H04),
as well as the Kafka immutability and DLT governance questions (H09, H13).
This concentration reveals a recurring archetype in which external compliance
obligations materialized into architectural hypotheses — a pattern discussed in
Section~\ref{sec:group-b-patterns} as the \textit{Regulatory Trigger Hypothesis}.

\subsection{Uncertainty and Impact Profile}
\label{subsec:uncertainty-impact}

Nine of the thirteen hypotheses carry a \textit{Very High} or \textit{High}
impact rating, signalling that the teams were not registering peripheral concerns
but uncertainties with potentially significant consequences for schedule, rework,
or system quality. Uncertainty ratings are more evenly distributed: five
hypotheses are rated \textit{Very Low} (the team has strong evidence but the
hypothesis is tracked for governance reasons), four are \textit{Medium}, one
is \textit{High}, and one \textit{Very High}.

The combination of Very Low uncertainty and Very High impact (H06, H07, H08,
H13) is particularly instructive. In ArchHypo's model, this quadrant maps to
hypotheses that \textit{can be postponed} — the team already has sufficient
evidence to act, so the immediate risk is low — yet the downstream consequences
of acting incorrectly are severe. Hypo-Stage's prioritization dashboard surfaced
these hypotheses in the "Can Postpone" segment of its focus panel, enabling
teams to consciously defer them without losing visibility. This represents
a practical instance of the \textit{Postpone Uncertain Decisions} pattern
described in the ArchHypo pattern literature \citep{silva2024patterns}.

\subsection{Technical Plans in Use}
\label{subsec:technical-plans}

Seven of the thirteen hypotheses had at least one associated technical plan at
the end of the observation period. Together these seven hypotheses carried thirteen
individual technical planning entries. The complete catalog of plan titles and
strategies is reproduced in Section~\ref{sec:app-technical-plans} of Appendix~B.
Table~\ref{tab:plan-strategies} summarizes
the strategies used.

\begin{table}[htbp]
\caption{Technical plan strategies observed across the hypothesis data.}
\label{tab:plan-strategies}
\centering\small
\begin{tabular}{lcc}
\toprule
\textbf{Strategy} & \textbf{Plans} & \textbf{Hypotheses using it} \\
\midrule
Guideline Implementation & 3 & 3 \\
Spike (Architectural Spike) & 3 & 2 \\
Modularization / Flexibility & 2 & 1 \\
Experiment & 2 & 2 \\
Tracer Bullet & 1 & 1 \\
Software Analytics (Trigger) & 1 & 1 \\
Other (implementation) & 1 & 1 \\
\midrule
\textbf{Total} & \textbf{13} & — \\
\bottomrule
\end{tabular}
\end{table}

Guideline Implementation and Architectural Spikes were the most frequently
chosen technical plans, consistent with findings from the PropSEP/Jetsoft evaluation
\citep{silva2025archhypo}. The Tracer Bullet plan appeared once for hypothesis H10,
where the Platform~Team planned a controlled injection of transient failures into
a real Kafka consumer path to measure ordering and latency characteristics. This
application illustrates how the Tracer Bullet differs from an Architectural Spike. Rather than
an isolated experiment, the Tracer Bullet plan exercises the full production consumer path
under representative conditions \citep{silva2024agilepatterns}.

Six hypotheses had no technical plan at the end of the study. Three of these
(H01, H04, H02) were registered but not yet acted upon; two (H07, H08) had
been validated through planning actions that predated the observation window and
whose plans were registered retrospectively; and one (H05) was under active
discovery but the plan document was still being prepared.

\subsection{Lifecycle and Status}
\label{subsec:lifecycle-status}

At the end of the observation period, ten hypotheses had \textit{Open} status,
one was \textit{In Review} (H05), and two were \textit{Validated} (H07 and H08,
both in the document storage domain). No hypothesis was invalidated or discarded
during the study window. The predominance of the Open status is consistent with
the short duration of the study (approximately nine calendar weeks, five two-week agile sprints spanned by hypothesis registrations) and with the
nature of the uncertainties: most concern mid- to long-horizon architectural
concerns that teams do not expect to resolve within a single sprint cycle.

% ─────────────────────────────────────────────────────────────
\section{RQ3 — Perceived Benefits and Challenges}
\label{sec:rq3-results}
\label{sec:results-rq3}
% ─────────────────────────────────────────────────────────────

\subsection{Perceived Benefits}
\label{subsec:benefits}

Active listening sessions with all four interviewed participants (the two
technical leaders from each team) converged on three recurring themes regarding
the perceived value of using Hypo-Stage.

\paragraph{Externalization of tacit architectural knowledge:}
Both teams described the act of registering a hypothesis as valuable in itself.
Participants noted that seeing a hypothesis written down with a clear statement,
an explicit uncertainty rating, an impact rating, and a related quality attribute
made it easier to share the concern with other colleagues who had not been involved
in the original discussion.

\paragraph{Prioritization support:}
The \textit{Where to Focus} panel of the Hypo-Stage dashboard, distinguishing
hypotheses that \textit{Need Attention} (high uncertainty and high impact) from
those that \textit{Can Postpone} (low impact), was described by both technical
leaders as a useful input in architectural conversations. One participant noted that
the "Can Postpone" category in particular gave them a way to explain
to other engineers outside of the team why a known concern was being deliberately deferred rather
than being ignored.

\paragraph{Collective memory for architectural decisions:}
Both teams observed that the hypothesis backlog served a documentation function
that was distinct from Architecture Decision Records (ADRs) or design documents.
Whereas ADRs capture decisions after they are made, an architectural hypothesis
captures an open question before a decision is available. Participants described
this as filling a gap in their current toolchain: they had good mechanisms for
recording what was decided but poor mechanisms for recording what was still
unknown.
Chapter~\ref{chap:conclusion} revisits this distinction when discussing implications for practice.

\subsection{Perceived Challenges}
\label{subsec:challenges}

Three challenges were consistently raised across participants.

\paragraph{Negative-formulation learning curve:}
The convention of writing hypotheses as falsifiable statements, phrased in a
form that specifies under which evidence the team would conclude the hypothesis
is incorrect, was described by all four participants as non-intuitive. Engineers
are accustomed to stating what they believe to be true. Inverting that belief
into a falsifiable form that describes what would prove it false requires a
deliberate cognitive effort that is easy to skip under time pressure. This finding
replicates the hard learning curve reported in both prior ArchHypo evaluations
\citep{silva2024tse, silva2024ieeesw}.

\paragraph{Registration time cost:}
Both technical leaders described the difficulty of dedicating the focused time
required to register a hypothesis properly. Fast-moving architectural micro-decisions, 
usually resolved within a day or less, were rarely registered because, by the time a
hypothesis could be written, the question had already been answered informally.
This creates a systematic under-representation of short-lifecycle uncertainties
in the hypothesis backlog, which may be consequential when the same question
re-emerges in a future context.

\paragraph{Absence of team-visible ceremonies:}
Neither team had established a recurring ceremony, such as an architectural
review or a hypothesis triage meeting, in which Hypo-Stage records were reviewed
as a team. Tool usage remained largely individual during the observation period.

% ─────────────────────────────────────────────────────────────
\section{RQ4 — Refinements to Hypo-Stage}
\label{sec:rq4-results}
\label{sec:results-rq4}
\label{sec:refinements}
% ─────────────────────────────────────────────────────────────

This section presents the refinements implemented in Hypo-Stage from commit
\texttt{ceee509} (the artifact boundary established by \citealt{correa2025}) to
commit \texttt{3a53d4e}. Each refinement was traced
to a concrete observation from the study and linked to the research question it
primarily addresses. All refinements address RQ4 by definition; selected
refinements also related to RQ1 (adoption ease) and RQ3 (usability and learning
curve).

\subsection{Backstage Integration and Catalog Surfacing}
\label{subsec:refinement-backstage}

\textbf{Observation:} During the training session and the first week of
observation, both teams navigated to Hypo-Stage via a standalone entry point
that was separate from their regular Backstage workflow. Engineers who were not
technical leaders rarely revisited the tool after the initial session.

\textbf{Refinement:} Hypo-Stage was wired as a first-class Backstage experience,
with routable pages for home, hypothesis creation, detail view, and edit flows.
Crucially, a catalog component tab (\texttt{EntityHypothesesTab}), shown in Figure~\ref{fig:entity-hypotheses-tab}, was introduced,
surfacing all hypotheses linked to the component an engineer is currently viewing
in the Backstage catalog.

\textbf{Impact:} Reducing navigation friction lowered the perceived cost of
checking the hypothesis backlog while already working with a service's catalog
entry.

\subsection{Backend API and Filtering}
\label{subsec:refinement-api}

\textbf{Observation:} Engineers with cross-team or platform-wide responsibilities
needed to filter hypotheses by team and by specific service. The initial
version of Hypo-Stage returned all hypotheses without server-side filtering,
which degraded usability as the hypothesis data grew.

\textbf{Refinements:} Three new API capabilities were implemented:
\begin{itemize}
  \item \texttt{GET /hypotheses} extended with \texttt{entityRef} and \texttt{team}
    query parameters;
  \item \texttt{GET /hypotheses/referenced-entity-refs} to populate the component
    filter without listing the entire catalog;
  \item \texttt{GET /hypotheses/teams} deriving team names from catalog metadata
    for components referenced by hypotheses.
\end{itemize}

\textbf{Impact:} On the Hypo-Stage home page, these capabilities enabled filtering
of the hypothesis list by team and by service (the Backstage catalog entity linked
to a hypothesis), so engineers with cross-team or platform-wide responsibilities
could narrow the backlog without loading the full unfiltered dataset.

\subsection{Prioritization Dashboard}
\label{subsec:refinement-dashboard}

\textbf{Observation:} The technical leaders expressed the desire to quickly
identify which hypotheses deserved immediate attention and which could safely be
deferred. Manually scanning the full hypothesis list imposed an unnecessary
cognitive burden, especially when the hypothesis data grew.

\textbf{Refinement:} A \texttt{HypothesisListDashboard} component was introduced
above the hypothesis table (Figure~\ref{fig:prioritization-dashboard}), presenting:
\begin{itemize}
  \item Overview cards with total count, validated count, and hypotheses created
    in the last thirty days;
  \item A \textit{Where to Focus} section distinguishing \textit{Need Attention}
    (high uncertainty and high impact) from \textit{Can Postpone} (low impact),
    with explanatory subtext aligned with ArchHypo's prioritization model
    \citep{silva2025archhypo};
  \item Distribution panels for uncertainty and impact Likert values.
\end{itemize}
The dashboard respects entity, team, and component scoping.

\textbf{Impact (RQ3).} The "Can Postpone" segment was specifically cited by
participants as resolving a communication challenge: it provided a principled
vocabulary for explaining deliberate deferral to non-technical stakeholders.

\begin{figure}[ht]
  \centering
  \includegraphics[width=0.85\textwidth]{prioritization-dashboard}
  \caption{Prioritization Dashboard in Hypo-Stage: the focus panel surfaces hypotheses
    according to the urgency signals derived from their uncertainty--impact combination,
    operationalizing the Postpone Uncertain Decisions pattern from the ArchHypo pattern literature \citep{silva2024patterns}.}
  \label{fig:prioritization-dashboard}
\end{figure}

\subsection{Hypothesis Lifecycle Consistency}
\label{subsec:refinement-lifecycle}

\textbf{Observation:} In the initial implementation, create and update operations
were not wrapped in database transactions, which created a risk of partial writes
(hypothesis row updated without the corresponding event entry, or vice versa).
Artifact analysis revealed two records with missing event entries, suggesting that
this race condition had occurred in practice.

\textbf{Refinements:} Create and update paths were refactored to use explicit
database transactions. Delete was extended to cascade to technical planning rows
and events, with clear API semantics and a confirmation dialog that requires the
exact hypothesis statement before proceeding (paste-enabled to avoid penalizing
users with long statements). Success and error toast notifications were added
for destructive operations.

\subsection{Technical-Planning-Driven Assessment Updates}
\label{subsec:refinement-tp-assessment}

\textbf{Observation:} Engineers reported that after completing a technical plan
action, such as an architectural spike, they updated the hypothesis's uncertainty
rating separately from the plan record. This two-step process was error-prone
and led to hypotheses whose assessment was stale relative to their technical
plans.

\textbf{Refinements:} Technical planning create and update operations were
extended to accept optional \texttt{uncertainty} and \texttt{impact} fields.
When present, the service updates the parent hypothesis's Likert values in the
same transaction. The event model was refined to distinguish
\texttt{TECHNICAL\_PLANNING\_CREATE} and \texttt{TECHNICAL\_PLANNING\_UPDATE}
events from generic hypothesis updates, preserving a traceable audit trail of
which assessment changes were driven by plan actions.

\subsection{Evolution Chart Improvements}
\label{subsec:refinement-chart}

\textbf{Observation:} The initial uncertainty/impact evolution chart used opaque
numeric y-axis ticks that participants could not interpret without consulting the
Likert scale documentation separately. Multiple events on the same calendar day
produced ambiguous x-axis labels.

\textbf{Refinements:} The chart was updated to render verbal y-axis ticks
(Very Low through Very High). Event source was made visually distinguishable:
circles for manual or creation-driven points, squares for technical-planning-driven
points. A legend explaining the marker shapes was added. When multiple events
occur on the same calendar day, the x-axis uses date and time to avoid ambiguity.
Tooltips were improved to label technical-planning-driven assessment moves with
the plan number when identifiable.

\begin{figure}[ht]
  \centering
  \includegraphics[width=0.8\textwidth]{evolution-chart}
  \caption{Evolution Chart in Hypo-Stage: tracking uncertainty and impact rating
    changes across successive assessments for a registered hypothesis. Each point
    corresponds to a saved assessment; the chart makes visible whether technical
    plan actions are reducing uncertainty or impact over time.}
  \label{fig:evolution-chart}
\end{figure}

\subsection{Form Usability Improvements}
\label{subsec:refinement-forms}

\textbf{Observation:} Participants found the hypothesis creation form cognitively
demanding, particularly the uncertainty and impact selectors, which displayed
only numeric values, and the quality attribute field, which accepted free text
without suggesting valid values.

\textbf{Refinements:} Uncertainty and impact selectors were redesigned as Likert
star pickers showing both a numeric button (1–5) and an adjacent text label
(Very Low through Very High) with disambiguating icons. Quality attributes were
replaced by a multi-select autocomplete showing attribute name and description
in the dropdown. Entity references for linking hypotheses to catalog components
were given a catalog-backed autocomplete. Documentation link fields were given an
improved list editor with add/remove and basic URL validation. 
Figure~\ref{fig:hypothesis-form} shows the improved hypothesis creation form.

\begin{figure}[ht]
  \centering
  \includegraphics[width=0.75\textwidth]{hypothesis-form}
  \caption{Hypothesis creation form in Hypo-Stage after the usability-oriented form
    refinements in this section: inline contextual guidance adjacent to the
    Statement field scaffolds the falsifiability inversion required by ArchHypo,
    reducing the cognitive barrier at the moment of registration.}
  \label{fig:hypothesis-form}
\end{figure}

\subsection{Summary: Refinements Mapped to Observations}
\label{subsec:refinement-summary}

Table~\ref{tab:refinements-map} summarizes the nine categories of refinements,
the primary observation that motivated each, and the research question each
addresses most directly.

\begin{table}[htbp]
\caption{Refinements to Hypo-Stage mapped to motivating observations and
  research questions. RQ columns indicate primary (P) or secondary (S) linkage.}
\label{tab:refinements-map}
\centering\small
\setlength{\tabcolsep}{4pt}
\begin{tabular}{p{3.8cm}p{4.5cm}ccc}
\toprule
\textbf{Refinement} & \textbf{Motivating observation} &
\textbf{RQ1} & \textbf{RQ4} & \textbf{RQ3} \\
\midrule
Backstage catalog tab
  & Infrequent revisit by non-leaders outside native workflow
  & S & P & S \\
Backend API filtering
  & Cross-team users needed scoped views as hypothesis data grew
  & S & P & — \\
Prioritization dashboard
  & Manual scanning of full backlog was cognitively costly
  & — & P & S \\
Lifecycle transaction consistency
  & Partial writes produced stale event records
  & — & P & — \\
TP-driven assessment updates
  & Assessment and plan records drifted apart post-spike
  & — & P & S \\
Evolution chart improvements
  & Numeric y-axis and same-day collisions caused confusion
  & — & P & S \\
Form usability improvements
  & Numeric-only selectors; free-text QA field caused errors
  & — & P & P \\
Safer deletion UX
  & Long statements made confirmation hostile without paste
  & — & P & — \\
Toast notifications
  & No feedback after destructive operations caused re-submission
  & — & P & — \\
\bottomrule
\end{tabular}
\end{table}

% ─────────────────────────────────────────────────────────────
\section{Emerging Patterns from the Case Study}
\label{sec:emerging-patterns}
% ─────────────────────────────────────────────────────────────

Beyond the direct answers to RQ1–RQ4, thematic analysis of observation notes,
active listening transcripts, and the combined hypothesis data produced two
families of patterns. Group~A describes team-behavior conditions that shape
whether hypothesis engineering becomes a sustainable practice. Group~B describes 
emerging patterns in the content of the architectural hypotheses
themselves. These two families address orthogonal dimensions: Group~A asks
\textit{when} hypothesis engineering works; Group~B asks \textit{what shape}
hypotheses take in practice. Table~\ref{tab:emerging-patterns} at the end of
this section summarizes the nine patterns across both groups.

\begin{quote}
\textbf{Scope and statistical caveat:}
The patterns described in this section were identified through thematic analysis
of thirteen architectural hypotheses produced by two software engineering teams
within a single organizational context over an observation period of approximately nine calendar weeks (five two-week agile sprints).
This dataset is sufficient to motivate conceptual characterization and to propose
initial vocabulary for the registered hypothesis, but it does not support statistical
claims about frequency, generalizability, or prevalence.
The patterns in Group A and Group B should therefore be read as
\emph{emerging propositions} --- grounded in real empirical data and theoretically
coherent --- but requiring replication across additional teams, organizations, and
domains before they can be treated as consolidated findings.
Future work directions for each pattern are discussed in Section~\ref{sec:future-work}.
\end{quote}

% ─────────────────────────────────────────────────────────────
\subsection{Group A — Team Behavior Patterns for Hypothesis Engineering
            Adoption}
\label{sec:group-a-patterns}
% ─────────────────────────────────────────────────────────────

The four Group~A patterns emerge by cross-referencing RQ1 and RQ3 findings. They
describe organizational and individual-level conditions that the study observed
to facilitate or obstruct the adoption of hypothesis engineering as a team
practice. These patterns extend the existing ArchHypo patterns observed by
\citet{silva2024patterns, silva2023initialpatterns} by addressing the social and
operational conditions in which ArchHypo mechanics operate in practice.

\emph{Note:} The four conditions described below are emergent rather than
confirmatory findings. Each is supported by convergent evidence from both
teams but they have not been tested in the broader population of software engineering
organizations. They are presented here as candidate patterns in descriptive
form; the derivation of actionable solutions from each pattern has been delegated to
a future-work research agenda in (Section~\ref{sec:future-work}), which outlines
the need for empirical validation before adoption as guidelines.

\subsubsection{Senior-Curated Hypothesis Backlog}
\label{subsubsec:pattern-a1}

\textbf{Problem:} In a team whose engineers operate at different scopes and
timescales, who is positioned to surface and sustain the registration of 
architectural hypotheses in mid- and long-terms?

\textbf{Observed evidence:} Twelve of the thirteen registered hypotheses were
authored by senior technical leaders — engineers who carried explicit
architectural responsibility and whose work extended beyond the current
sprint. Non-senior engineers did not sustain independent hypothesis registration
after the initial training episode. The hypotheses that were registered shared
a common trait: they concerned architectural uncertainties whose consequences would materialize
weeks or months in the future, making them invisible to typical sprint-scope
review practices.

\subsubsection{Fast-Resolution Blind Spot}
\label{subsubsec:pattern-a2}

\textbf{Problem:} Many architectural micro-decisions are made and resolved within
hours or a single day. Their speed makes them feel too transient to register, yet
their cumulative architectural effect may only become visible, as a drift or an
inconsistency, long after the original decision context is forgotten.

\textbf{Observed evidence:} Both technical leaders reported that the majority of
architectural decisions encountered day-to-day were resolved before a hypothesis
could be written. This speed made registration feel retrospectively unnecessary.
Yet both also described recurring situations where the same question re-surfaced
in a new context, requiring the team to reconstruct reasoning that had previously
been resolved informally and in a rapid pace.

\subsubsection{Registration as a Slack-Dependent Practice}
\label{subsubsec:pattern-a3}

\textbf{Problem:} Translating a tacit architectural concern into a properly
formulated ArchHypo hypothesis requires focused, uninterrupted attention.
Sprint delivery pressure makes this kind of attention scarce and unreliable as
a basis for sustained practice.

\textbf{Observed evidence:} Senior leaders explicitly described the difficulty
of "dedicating proper time to register all possible hypotheses they encounter
day-to-day." The observed set of thirteen hypotheses is known to
under-represent the total uncertainties that leaders were aware of. Registration
activity clustered around protected and restricted discussion time rather than being distributed
across normal sprint work.

\subsubsection{Negative-Formulation Barrier}
\label{subsubsec:pattern-a4}
\textbf{Problem:} The ArchHypo convention of writing hypotheses as falsifiable
statements represents a cognitive inversion of how engineers typically articulate their beliefs.
This inversion is a persistent adoption barrier even for practitioners who
understand the technique conceptually.

\textbf{Observed evidence:} All four interviewed participants mentioned that the
falsifiability requirement is non-intuitive and difficult to apply under time
pressure. This replicates findings reported in the PropSEP/Jetsoft evaluation
\citep{silva2025archhypo} and remains unresolved in the current version of
the technique's guidance materials.

% ─────────────────────────────────────────────────────────────
\subsection{Group B — Technical Architectural Patterns for Hypothesis
            Engineers}
\label{sec:group-b-patterns}
% ─────────────────────────────────────────────────────────────

The prior work in the ArchHypo research line has established the technique's
core concepts \citep{silva2025archhypo}, documented patterns for hypothesis
assessment, uncertainty management, and decision deferral
\citep{silva2024patterns, silva2023initialpatterns}, and reported on the
technique's learning curve and process-integration challenges
\citep{silva2024archhypo}. What this body of work has not yet addressed is
a question that only becomes answerable through empirical data on
hypotheses registered in practice:

\begin{quote}
\textit{What recurring shapes do architectural hypotheses take when teams apply
ArchHypo to real systems? And what structural properties of a hypothesis, its
source type, quality-attribute profile, and relation to its organizational
context, should an engineer recognize in order to write and manage it more
effectively?}
\end{quote}

This gap was also identified in calls for future work within the ArchHypo
pattern literature, which explicitly noted the need for "identifying and
documenting patterns for various types of hypotheses" and for "exploration of
common uncertainties related to specific quality attributes"
\citep{silva2024patterns}. The thirteen-hypothesis dataset produced by this
study constitutes the first empirical dataset that enables such identification.
The five patterns described below represent a \textbf{novel contribution}
of this work: they are grounded in real hypothesis data, not derived
from theoretical analysis, and they provide a vocabulary that 
engineers can use both to recognize familiar situations more quickly and to
communicate about hypothesis types across teams.

\emph{Statistical scope:} All five patterns were identified from the same
thirteen-hypothesis dataset. While each pattern maps to two or more
hypotheses in the dataset, the small dataset size means that no frequency claims are made. 
The contribution lies in naming and structuring recurring shapes so that future practitioners and researchers
have a vocabulary for recognition, not in establishing that any pattern
will appear with a particular frequency in other organizational contexts, nor
in prescribing solutions for each one. Solution-shaped guidance and
consolidation into the broader ArchHypo pattern language are delegated to the
future-work research agenda in (Section~\ref{sec:future-work}).

\subsubsection{Hypothesis Cluster}
\label{subsubsec:pattern-b1}

\textbf{Context:} A team is responsible for a shared platform component on which
multiple other service teams depend. An architectural question about that
component cannot be decomposed into a single hypothesis because the design space
has several interrelated dimensions, each of which involves a distinct trade-off.

\textbf{Problem:} Formulating a single hypothesis for a multi-dimensional
architectural concern either oversimplifies the problem (losing important
trade-off structure) or produces an unmanageably broad statement that is hard
to falsify.

\textbf{Observed evidence:} Five hypotheses (H09–H13) all concern the shared
Kafka streaming infrastructure and were registered on the same day, as the
Platform~Team's technical leader used Hypo-Stage to structure a design discussion.
The five hypotheses address: (i)~DLT audit semantics versus offset-based replay;
(ii)~synchronous versus asynchronous retry for transient consumer failures;
(iii)~failure-state placement inside versus outside the Kafka flow;
(iv)~error classification before retry; and (v)~payload immutability before
re-enqueueing. Each is independently falsifiable yet together they map the full
decision space for consumer-side failure handling in the shared streaming
platform.

\subsubsection{Regulatory Trigger Hypothesis}
\label{subsubsec:pattern-b2}
\textbf{Context:} A software system operates in a regulated domain such as financial
services, healthcare, or data privacy, where external compliance obligations
create architectural constraints whose implementation path is not fully known.

\textbf{Problem:} A compliance requirement imposes an architectural obligation
before the team has determined how to satisfy it. It is unclear whether the
existing architecture already satisfies the obligation, what changes would be
required if it does not, and what the cost of non-compliance would be.

\textbf{Observed evidence:} Two hypotheses in the dataset arise directly from
regulatory obligations: H04 (the GDPR, the General Data Protection Regulation, right-to-erasure requirement for binary
documents in the blob storage service) and H03 (organization-wide security
audit requirements for database log correlation across the platform framework).
In both cases the uncertainty is not about whether to comply — that is a legal
requirement that must be met — but about whether the current architecture \textit{can} comply and at
what technical cost.

\subsubsection{Context-Agnostic Restructuring Hypothesis}
\label{subsubsec:pattern-b3}
\textbf{Context:} A service currently handles multiple responsibilities such as domain
logic, external integrations, and identity management. The team believes the service could be restructured to serve
multiple domain contexts, but what a context-agnostic restructuring would entail---contract boundaries and migration path---remains unknown.

\textbf{Problem:} Deciding \textit{whether} to pursue context-agnostic restructuring is itself uncertain: the
team does not know what the correct contract boundary should be, which existing
flows can be migrated without business-logic changes, and what the downstream
impact on other teams would be. Committing to that restructuring without this information
risks a disruptive refactor that delivers less reuse than expected.

\textbf{Observed evidence:} H05 (credit-bureau context-agnostic restructuring) captures this
precisely. The service currently handles identity, bureau integration, and
confidential document flows using domain-specific entities tied to a single
business vertical. The hypothesis asks whether these integrations can be
decoupled into context-agnostic middleware accepting only each provider's minimum
required identifiers, reusable across domains without duplicated cost. At the
end of the observation period, H05 was the only hypothesis in \textit{In Review}
status, with a discovery technical plan mapping (i)~the dependency matrix of
current versus real dependencies per integration, (ii)~an agnostic interface
contract proposal, and (iii)~a migration plan identifying which flows can isolate
immediately.

\subsubsection{Load-Validated Hypothesis}
\label{subsubsec:pattern-b4}
\textbf{Context:} A team has proposed a specific algorithmic or implementation
approach to handle a high-volume data processing challenge. The team has already
conducted a capacity or volume test at scale, producing evidence that the
approach works. However, the hypothesis has not been formally closed because
the implementation is still in progress.

\textbf{Problem:} The hypothesis occupies the lowest uncertainty quadrant
(Very Low) while retaining the highest impact rating (Very High), reflecting a
situation where the team \textit{knows} the solution works but the architectural
commitment is still consequential. The hypothesis must be tracked until
implementation is complete, yet its characterization as "Open" may suggest
unresolved doubt that no longer exists.

\textbf{Observed evidence:} H06 (billing aggregation chunked keyset pagination)
exemplifies this. The team validated the approach — chunked keyset pagination
with isolated repository-level transactions — through a capacity test at five
times production scale (45 million rows total, 1.5 million rows per month).
Profiling confirmed that JPA stream processing caused first-level cache growth,
and that chunking scopes transactions to approximately ten thousand rows so the
garbage collector can reclaim memory, producing a healthy sawtooth heap pattern.
The hypothesis remains Open because the implementation is ongoing, but the
uncertainty has been effectively resolved.

\subsubsection{Shared Resource Governance Hypothesis}
\label{subsubsec:pattern-b5}
\textbf{Context:} A platform service manages a shared resource such as binary objects,
credentials, and configuration data, on which multiple product domains depend. New
business requirements demand that the resource be shared across domain boundaries
or transferred between owning services, but the governance rules for these
operations have not been established.

\textbf{Problem:} The team must decide how to model cross-domain access to the
shared resource — specifically, which service "owns" the resource at any given
point, who may read or write it, and under what conditions ownership can be
transferred. These questions have significant security, consistency, and usability
implications. Without explicit governance decisions, different product teams will
adopt incompatible access patterns, creating consistency violations that are
difficult to detect and remedy.

\textbf{Observed evidence:} H07 (blob ownership transfer across service domains)
and H08 (blob cross-client sharing governance) are two instances of this pattern
in the secure blob storage domain. H08 was driven by a concrete usage scenario
in which a customer-facing channel receives documents that are later routed to
support tooling; the support team needs read access without breaking the blob's ownership invariants. 
Both hypotheses were validated through a Guideline technical plan — an alignment session that produced a governance decision enabling
prioritization of the transfer and sharing implementation work.

\begin{table}[htbp]
\caption{Emerging patterns identified in the case study: Group~A (team-behavior
conditions for sustainable hypothesis engineering) and Group~B (structural
archetypes of architectural hypotheses). Solution-shaped guidance for each
pattern is deferred to the future-work research agenda
(Section~\ref{sec:future-work}).}
\label{tab:emerging-patterns}
\centering\small
\setlength{\tabcolsep}{4pt}
\begin{tabular}{clp{3.6cm}p{5.2cm}p{2.6cm}}
\toprule
\textbf{ID} & \textbf{Group} & \textbf{Pattern name} & \textbf{Signature} &
\textbf{Primary evidence} \\
\midrule
A1 & Team behavior & Senior-Curated Hypothesis Backlog
   & Mid-/long-horizon uncertainties stay visible because senior engineers
     author them
   & 12 of 13 hypotheses senior-authored \\
A2 & Team behavior & Fast-Resolution Blind Spot
   & Architectural micro-decisions resolve before registration, risking
     cumulative drift
   & Both technical leaders' testimony \\
A3 & Team behavior & Registration as a Slack-Dependent Practice
   & Well-formulated hypotheses require protected time, not sprint slack
   & Registration clustering at training and Kafka session \\
A4 & Team behavior & Negative-Formulation Barrier
   & Falsifiable inversion is persistently cognitively costly
   & All four interviewed participants \\
\midrule
B1 & Structural archetype & Hypothesis Cluster
   & Multi-dimensional platform concern decomposes into linked falsifiable
     hypotheses
   & H09--H13 (Kafka failure handling) \\
B2 & Structural archetype & Regulatory Trigger Hypothesis
   & External compliance obligation surfaces as tracked architectural
     uncertainty
   & H03, H04 \\
B3 & Structural archetype & Context-Agnostic Restructuring Hypothesis
   & Uncertainty about whether and how to restructure toward context-agnostic boundaries
   & H05 (credit-bureau context-agnostic restructuring) \\
B4 & Structural archetype & Load-Validated Hypothesis
   & Very-low uncertainty, very-high impact: solution empirically validated
     but implementation open
   & H06 (billing aggregation) \\
B5 & Structural archetype & Shared Resource Governance Hypothesis
   & Cross-domain access and ownership rules for a shared platform resource
   & H07, H08 \\
\bottomrule
\end{tabular}
\end{table}

% ─────────────────────────────────────────────────────────────
\section{Revisiting the Research Questions}
\label{sec:rq-revisit}
% ─────────────────────────────────────────────────────────────

This section concludes the analytical observations of the chapter by stating explicitly how
the preceding empirical and design material responds to the four research questions
formulated in Chapter~\ref{chap:introduction}. The intent is \emph{consolidative} where
no additional empirical claims are introduced. The aim is to make the logical
organization of the chapter legible relative to the research questions.

\subsection{Correspondence between research questions and chapter sections}

Table~\ref{tab:rq-chapter-map} summarizes where each research question is principally
addressed in this chapter, and where integrated interpretation (Group~A/B) extends
the principal sections.

\begin{table}[htbp]
\centering
\small
\setlength{\tabcolsep}{5pt}
\caption{Correspondence between the research questions (Chapter~\ref{chap:introduction})
  and the organization of evidence in this chapter.}
\label{tab:rq-chapter-map}
\begin{tabular}{p{0.12\linewidth} p{0.40\linewidth} p{0.40\linewidth}}
\toprule
\textbf{RQ} & \textbf{Principal analysis} & \textbf{Integrated interpretation} \\
\midrule
RQ1 & Section~\ref{sec:rq1-results} (tool use and adoption dynamics). &
Section~\ref{sec:group-a-patterns} (Group~A; read together with RQ3). \\
\midrule
RQ2 & Section~\ref{sec:rq2-results} (typological account of the data). &
Section~\ref{sec:group-b-patterns} (Group~B structural archetypes). \\
\midrule
RQ3 & Section~\ref{sec:rq3-results} (benefits and challenges). &
Section~\ref{sec:group-a-patterns} (Group~A; read together with RQ1). \\
\midrule
RQ4 & Section~\ref{sec:rq4-results} (nine refinement categories, observation-led). &
--- \\
\bottomrule
\end{tabular}
\end{table}

\subsection{Brief interpretive synthesis}

Through the preceding analysis presented in this chapter, the four research questions tell a single story about how far a
tool-mediated instantiation of ArchHypo can carry a team, and where the remaining barriers lie.

\textbf{RQ1 and RQ3 converge on an organizational, not a tool, problem:}
Both questions produce evidence that Hypo-Stage lowered surface friction, 
facilitating context-switching, filtering, traceability. Participants
recognized the intellectual value of hypothesis management. Still, the same two
questions also show that authorship stayed concentrated in senior engineers,
who registered hypotheses clustered around protected discussion time rather than
distributing across sprint work, and that the falsifiability inversion
remained cognitively costly. The Group~A patterns (A1--A4) name the
conditions under which those findings hold: adoption breadth is limited by the
organization structure including senior curation, protected cognitive time,
statement-level support, and ceremonial touchpoints. These are conditions that a tool alone
cannot supply.

\textbf{RQ2 and Group~B reveal a content gap in the existing ArchHypo
pattern language:} The prior pattern literature
\citep{silva2024patterns, silva2023initialpatterns} names mechanics of
assessment, deferral, and uncertainty management, but does not name the
\emph{kinds} of hypotheses engineers actually register. When applied in a
compliance-heavy, platform-dominant environment, the thirteen registered
hypotheses concentrate on regulatory triggers, shared-resource governance,
context-agnostic restructuring, and multi-dimensional platform concerns, which are categories the
existing patterns do not index. Group~B therefore complements the
mechanics-focused patterns rather than competing with them, and provides
the first empirically grounded vocabulary for the content of hypotheses
that emerge in practice.

\textbf{RQ4 closes the formative-evaluation loop:} The nine refinement
categories are not a separate implementation narrative. Each maps to a
specific observed friction surfaced by RQ1, RQ2, and RQ3 questions: the form-level and
ceremony-related refinements track A3 and A4; the component-scoped tab
tracks B1-shaped concerns about multi-dimensional platform uncertainty; the
prioritization dashboard and API filtering respond to the adoption
dynamics observed under RQ1. Their derivation from the same observational
data is itself evidence that the design-science loop of observation,
refinement, and re-observation operated on the Hypo-Stage tool during the study.

Taken together, the chapter supports a bounded claim: tool-mediated
hypothesis engineering succeeds at the integration and traceability layers
and produces a useful vocabulary for the content of architectural
hypotheses, but its adoption ceiling lies at the organizational layer. The
Group~A patterns are the natural hand-off to
Chapter~\ref{chap:conclusion}, where they are read as open empirical problems rather than as actionable guidelines.

%!TeX root=../thesis.tex

\chapter{Conclusion}
\label{ch:conclusion}
\label{chap:conclusion}

This chapter closes the empirical and analytical work of the dissertation.
Section~\ref{sec:conclusion-rq} summarizes the main contributions in relation
to each research question. Section~\ref{sec:implications} discusses
implications for practitioners and researchers.
Section~\ref{sec:future-work} outlines the future work agenda that this study
opens.

% ─────────────────────────────────────────────────────
\section{Summary in Relation to the Research Questions}
\label{sec:conclusion-rq}

\paragraph{RQ1 --- How do software engineering teams use Hypo-Stage and ArchHypo
when working on a scaled product?}

The study produced a detailed observational account of tool adoption in two
teams over four weeks of real sprint work. The most consistent finding is
that hypothesis registration was senior-led: twelve of the thirteen registered
hypotheses were authored by staff or principal engineers who held explicit
architectural responsibilities. Junior and mid-level engineers were aware of
the tool and trained on the technique, but did not register hypotheses
independently within the observation window. Registration activity was
concentrated in two bursts --- one coinciding with the training session, and a
second when the Platform Team used the tool to structure a live architectural
discussion on Kafka failure-handling strategy --- rather than being continuously
distributed across sprints.

The integration of Hypo-Stage with the Backstage Internal Developer Portal
was consistently cited as an enabler of adoption. The
\texttt{EntityHypothesesTab} --- which surfaces only hypotheses linked to the
service a developer is already viewing --- reduced the cognitive cost of
context-switching into hypothesis management and was specifically mentioned as
a meaningful design improvement over a standalone tool. However, embedding
the tool into team-visible ceremonies (retrospectives, planning sessions)
did not occur spontaneously within the observation window; both teams noted that
this would require an explicit cultural agreement.

\paragraph{RQ2 --- What types and recurring structural shapes of architectural
uncertainties do teams identify, record, and manage using Hypo-Stage?}

The study produced thirteen architectural hypotheses across two teams. Analysis
of their source types, quality attributes, uncertainty and impact profiles, and
technical plan strategies characterizes the landscape of architectural
uncertainty that Hypo-Stage users made explicit (Section~\ref{sec:rq2-results}).
Requirements-sourced hypotheses dominated the Product Team (five of six), while
Solution-sourced hypotheses dominated the Platform Team's Kafka cluster (four of
five). The dominant quality attributes --- Reliability and Auditability ---
reflect the teams' real operational context: a scaled financial services product
operating under both internal security audit and privacy regulation obligations.

Beyond this characterization, five recurring structural archetypes were
identified from the data through thematic analysis: \textit{Hypothesis Cluster},
\textit{Regulatory Trigger Hypothesis}, \textit{Decoupling Seam Hypothesis},
\textit{Load-Validated Hypothesis}, and \textit{Shared Resource Governance
Hypothesis} (Section~\ref{sec:group-b-patterns}). These archetypes constitute the first empirically
grounded answer to the question of what shapes architectural hypotheses take when
ArchHypo is applied in a real industrial context, and they directly address a gap
identified as future work in \citet{silva2024patterns}. Readers should note
that these are \emph{emerging patterns}, derived from a thirteen-hypothesis
dataset; they require replication in additional organizational contexts before
they can be treated as consolidated contributions to the ArchHypo pattern
language (see Section~\ref{sec:future-work}).

\paragraph{RQ3 --- What benefits and challenges do team members perceive in
using Hypo-Stage?}

Participants across both teams reported that making architectural uncertainties
explicit and visible through Hypo-Stage had a positive effect on the quality of
their architectural discussions. The most salient benefit was the formalization
of tacit knowledge: by writing a hypothesis, engineers converted an informal
concern --- shared informally in chat or resolved verbally in a meeting --- into a
traceable, prioritized artifact with an explicit rationale. The Backstage
catalog integration was perceived as reducing cognitive overhead.

The dominant challenge remained the learning curve of the ArchHypo technique
itself. All four participants in the active listening sessions identified the
falsifiability requirement of hypothesis statements as non-intuitive under
delivery pressure, replicating findings from the PropSEP/Jetsoft evaluation
\citep{silva2025archhypo}. Sprint delivery pressure was named by both technical
leaders as the principal constraint on registration activity: formulating a
well-structured hypothesis requires focused, uninterrupted time that is scarce
in a product sprint. The team-behavior patterns in Group~A (Section~\ref{sec:group-a-patterns})
condense these findings into four actionable organizational conditions ---
\textit{Senior-Curated Hypothesis Backlog} (A1), \textit{Fast-Resolution Blind Spot} (A2),
\textit{Registration as a Slack-Dependent Practice} (A3), and
\textit{Negative-Formulation Barrier} (A4) --- that jointly explain why adoption is
sustainable only when specific organizational enablers are in place. Like the
Group~B archetypes, these are emerging propositions and require further
validation.

\paragraph{RQ4 --- Which refinements to Hypo-Stage and its usage guidelines
emerge from the empirical study?}

Nine categories of empirically grounded refinements were implemented during the
observation period and documented in Section~\ref{sec:refinements}, each traced to a concrete
observation and a commit in the Hypo-Stage repository. The refinements span
four layers of the tool: Backstage integration and catalog surfacing (making
hypotheses discoverable without navigation); backend API and filtering (enabling
team-scoped and status-filtered views); the prioritization dashboard
(introducing focus panels and urgency signals derived from uncertainty and impact
combinations); and form usability (adding inline ArchHypo guidance to reduce the
vocabulary barrier at the moment of registration).

The most structurally significant refinement is the evolution chart introduced
in Section~\ref{subsec:refinement-chart} (Figure~\ref{fig:evolution-chart}), which plots the shift in a
hypothesis's uncertainty and impact ratings over successive assessments. This
feature directly addresses the absence of assessment history in the initial
Hypo-Stage version \citep{correa2025} and enables teams to see whether their
technical plan actions are having the intended effect of reducing uncertainty
or impact. Together, the nine refinement categories operationalize a cycle of
\emph{observe $\rightarrow$ refine $\rightarrow$ redeploy $\rightarrow$ observe} that is consistent with the
action-research element of the research design.

% ─────────────────────────────────────────────────────
\section{Implications for Practice and Research}
\label{sec:implications}

\paragraph{For practitioners adopting hypothesis-engineering tools.}

The study demonstrates that tool integration within the developer's daily
environment is a meaningful enabler of adoption. Deploying Hypo-Stage as a
Backstage plugin --- where engineers already manage services, documentation, and
CI/CD pipelines --- removed a class of adoption friction that a standalone tool
would have created. This suggests a general principle: hypothesis-engineering
tools are more likely to be adopted when they are delivered as extensions to
an existing workflow rather than as a new workflow to be added.

The finding that registration remains senior-led, even with tool support, has
practical implications for organizations intending to distribute architectural
decision-making. Tool availability is a necessary but not sufficient condition
for broad participation. Organizations should combine tool deployment with
explicit protected time for architectural hygiene (see Group~A pattern A3) and
with statement scaffolding that lowers the cognitive barrier to formulating
falsifiable hypotheses (see Group~A pattern A4 and Hypo-Stage form improvements
in Section~\ref{subsec:refinement-forms}). The Group~B archetypes can serve immediately as a
recognition vocabulary for experienced hypothesis engineers encountering
familiar situations in new projects.

\paragraph{For researchers studying tool-mediated architectural practices.}

This study is the first evaluation of ArchHypo in a tool-mediated setting and
the first to produce a dataset of architectural hypotheses large enough to
support structural classification. Several methodological observations may
benefit future work. First, the formative evaluation design --- in which
refinements are implemented during the observation period --- provides richer
insight into \emph{why} difficulties arise and \emph{which} changes address
them than a retrospective evaluation would. This design is well-suited to tool
artifacts whose value proposition is tightly coupled to usability and workflow
fit. Second, the dataset of thirteen hypotheses, while small, demonstrates that
even a short observation window in a naturalistic setting produces enough
artifacts to support initial pattern identification. Larger datasets will be
needed to test frequency and generalizability.

The observation that hypothesis registration is senior-led is consistent with
findings from both prior ArchHypo evaluations \citep{silva2024ieeesw,
silva2024tse} and with the broader literature on tacit knowledge in architectural
decision-making \citep{souza2019}. This convergence across different settings,
time periods, and methods provides some basis for treating senior-anchor
dependency as a robust characteristic of early ArchHypo adoption rather than an
artifact of any single study.

% ─────────────────────────────────────────────────────
\section{Limitations and Future Work}
\label{sec:future-work}

\subsection*{Study Limitations}

\paragraph{Single organizational context and Backstage dependency.}
The study was conducted within one technology organization that already operated
a mature Backstage deployment. The enabling condition of Backstage availability
means that findings about adoption friction may not transfer directly to
organizations without an Internal Developer Portal. The organizational culture,
team size, and product complexity of the study site may differ substantially from
other contexts in which researchers or practitioners seek to apply these findings.

\paragraph{Small hypothesis dataset and short observation window.}
Thirteen hypotheses produced over four weeks by two teams is a limited dataset.
The Group~A and Group~B patterns are derived from this dataset and must be treated
as emerging propositions rather than consolidated findings. The four-week window
is also too short to observe full hypothesis lifecycle completion (validation or
invalidation); eleven of the thirteen hypotheses remained in Open status at the
end of the period. A longer observation window would be needed to study how
teams manage hypotheses through to resolution and whether the tool adequately
supports retrospective learning after a hypothesis closes.

\paragraph{Researcher dual role.}
The researcher's dual role as organizational employee and academic investigator
is a known source of potential bias. Although strategies were employed to
maintain role separation (Section~\ref{sec:researcher-role}), the researcher's presence as a tool
author and subject-matter expert may have influenced how participants described
their experience of the technique, introducing social desirability effects.

\paragraph{No comparison condition.}
The study does not include a control group or comparison condition against teams
using ArchHypo without Hypo-Stage, or teams not using any structured
uncertainty-management approach. This limits the ability to attribute observed
benefits specifically to the tool rather than to the introduction of a structured
technique in general.

\subsection*{Future Work Agenda}

\paragraph{Replication and longitudinal extension.}
The most immediate priority is replication of this study in additional
organizational contexts --- including organizations without an existing Backstage
deployment --- to test whether the adoption patterns, tool refinements, and
emerging archetypes are context-specific or more broadly applicable.
A longitudinal design spanning multiple development cycles would allow the
study of hypothesis lifecycle completion, assessment evolution, and whether
teams internalize the technique over time to the point where senior-anchor
dependency diminishes.

\paragraph{Consolidating the Group~B archetypes into the ArchHypo pattern language.}
The five technical architectural hypothesis archetypes (B1--B5) were identified
from a single dataset and are explicitly preliminary. Consolidating them
requires: (a)~attempting to identify or refute each archetype in new hypothesis
datasets from different organizations and domains; (b)~subjecting each pattern
to the structured review process used for the existing ArchHypo pattern language
\citep{silva2024patterns}; and (c)~expanding the catalog to cover quality
attributes not represented in this study (e.g., security-only hypotheses,
performance-at-scale hypotheses in non-financial contexts).

\paragraph{Validating the Group~A team-behavior patterns.}
The four organizational conditions in Group~A (A1--A4) are equally preliminary.
Future work should design intervention studies that test whether addressing
these conditions --- for instance, by introducing protected architectural hygiene
sessions (A3) or embedding statement templates in the hypothesis form (A4) ---
measurably increases registration breadth beyond senior engineers, reduces
the formulation time reported by participants, or improves hypothesis quality
as judged by the teams themselves.

\paragraph{Broadening participation beyond senior engineers.}
A central open question is whether and how hypothesis registration can become a
practice accessible to engineers at all levels, not only to those who already
hold explicit architectural responsibilities. Research directions include:
guided hypothesis templates that lower the formulation barrier (already initiated
in the form usability refinements); asynchronous architectural review workflows
that do not require simultaneous availability of a senior engineer; and
gamified or social nudge mechanisms that encourage mid-level engineers to
surface uncertainties they currently resolve tacitly.

\paragraph{Integration with Architecture Decision Records.}
Hypo-Stage manages the pre-decision lifecycle of architectural uncertainty.
Once a hypothesis is validated or invalidated, the resulting decision is
currently not captured in any structured artifact within the tool. Future
work should explore the integration of Hypo-Stage with ADR tools such as
Log4brains, so that the evidence chain from hypothesis registration through
technical plan execution to final architectural decision is preserved as a
single traceable thread.

\paragraph{Quantitative assessment of tool impact.}
This study does not include any quantitative measure of the impact of
Hypo-Stage on decision quality, architecture drift, or delivery velocity.
Future work using a quasi-experimental or interrupted time-series design could
compare architectural decision quality or technical debt accumulation between
periods with and without active hypothesis management, providing a more
rigorous basis for claims about the tool's organizational value.

\paragraph{Extending Hypo-Stage to multi-team hypothesis governance.}
The present evaluation involved two co-located teams within the same
organization. As organizations scale, hypotheses registered by one team
frequently have cross-team implications --- this was already visible in the blob
governance hypotheses (H07, H08) and the platform Kafka cluster (H09--H13).
Future work should explore cross-team hypothesis visibility, shared ownership
models, and notification mechanisms that make cross-boundary hypotheses
discoverable to the teams they affect, without creating coordination overhead
that undermines team autonomy.



%%%%%%%%%%%%%%%%%%%%%%%%%%%% APPENDICES %%%%%%%%%%%%%%%%%%%%%%%%%%%%%%%%%%%%%%%%

% Appendices = your own supplementary content (questionnaires, protocols, etc.)
\cleardoublepage
\pagestyle{appendix}
\appendix
\addappheadtotoc

%!TeX root=../thesis.tex

\chapter{Research Instruments}
\label{apx:instruments}

This appendix reproduces the instruments and communications used during data
collection. Section~\ref{sec:listening-guide} contains the open-ended active
listening guide used in the end-of-period sessions.
Section~\ref{sec:informed-consent} reproduces the informed consent framework
communicated to participants. Section~\ref{sec:training-comm} reproduces the
training communication sent to both teams at the start of the observation period.

% ─────────────────────────────────────────────────────
\section{Open-Ended Active Listening Guide}
\label{sec:listening-guide}

\textbf{Note on instrument design.} The study did not employ a fixed
questionnaire. As described in Section~\ref{subsec:active-listening}, data collection followed an
open-ended active listening format consistent with the semi-structured
interview approach recommended by \citet{robson2016} for flexible designs.
The guide below provided a thematic scaffold for conversations; the actual
order and phrasing of questions varied with the natural flow of each session.
Sessions were conducted individually with the senior technical leader of each
team at the end of the observation period.

\subsection*{Thematic Areas and Guiding Questions}

\textbf{Area 1 --- Overall experience of Hypo-Stage.}
\begin{itemize}
  \item Can you walk me through the moments when you actually opened Hypo-Stage
    during the observation period?
  \item What made you decide to register a hypothesis at those particular
    moments?
  \item Were there situations where you thought about registering a hypothesis
    but ended up not doing so? What stopped you?
\end{itemize}

\textbf{Area 2 --- Perceived usefulness of the technique and the tool.}
\begin{itemize}
  \item Has using Hypo-Stage changed the way you think about or discuss
    architectural decisions with your team?
  \item Are there specific decisions you handled differently because you had a
    hypothesis registered?
  \item What aspect of the tool did you find most useful in practice?
  \item What aspect of the tool was most frustrating or felt like friction?
\end{itemize}

\textbf{Area 3 --- Learning curve and technique understanding.}
\begin{itemize}
  \item How confident do you feel formulating a hypothesis statement that is
    falsifiable? Did that feel natural or forced?
  \item How did your team members react to the concept of an architectural
    hypothesis? Do you think the concept would be accessible to a mid-level
    engineer without facilitation?
  \item What would have made the onboarding experience easier?
\end{itemize}

\textbf{Area 4 --- Workflow integration.}
\begin{itemize}
  \item How well did Hypo-Stage fit into your existing sprint ceremonies and
    day-to-day tooling?
  \item Did the Backstage integration help or was it irrelevant to adoption?
  \item If you were to continue using the tool after this study, what would
    need to change --- either in the tool or in your team's process --- for it to
    be sustainable?
\end{itemize}

\textbf{Area 5 --- Future outlook.}
\begin{itemize}
  \item If you had to name one architectural uncertainty that you are currently
    managing tacitly that you \emph{wish} you had registered in Hypo-Stage,
    what would it be?
  \item Would you recommend Hypo-Stage to another team in the organization?
    Under what conditions?
\end{itemize}

% ─────────────────────────────────────────────────────
\section{Informed Consent Framework}
\label{sec:informed-consent}

No formal written consent form was required for this study, as the research
involved no sensitive personal data and participation was entirely voluntary
(Section~\ref{sec:ethical-considerations}). The following points were communicated verbally and in writing
to all participants prior to the observation period and were confirmed by
engineering leadership in the organization:

\begin{enumerate}
  \item \textbf{Voluntary participation.} All engineers were free to use
    Hypo-Stage as much or as little as they chose. Participation in the
    active listening sessions was optional and had no bearing on the
    researcher's organizational role.
  \item \textbf{Anonymization.} The name of the organization, the names of
    all participants, and the names of internal software products and
    components would be anonymized in all research outputs. Backstage entity
    references in hypothesis data would be replaced with generic labels
    (see Appendix~\ref{app:hypotheses-data}, Table~\ref{tab:entity-legend}).
  \item \textbf{Data use.} Observations, artifact records, and session notes
    would be used exclusively for academic research purposes and would not be
    shared with organizational management in any individually identifiable form.
  \item \textbf{Right to withdraw.} Any participant could withdraw previously
    provided data or discontinue participation at any point without consequence.
  \item \textbf{Research knowledge and consent of leadership.} The study
    was conducted with the knowledge and informal consent of the relevant
    engineering leadership in the organization, who confirmed that participation
    was compatible with the teams' work obligations.
\end{enumerate}

% ─────────────────────────────────────────────────────
\section{Training Communication}
\label{sec:training-comm}
\label{apx:training-communication}

The following text reproduces, in anonymized and lightly edited form, the
written communication sent to both teams prior to the training session that
inaugurated the observation period. The purpose of this communication was to
set expectations, introduce the research context, and invite voluntary
participation.

\begin{quote}
\textbf{Subject: Architectural Hypothesis Management with Hypo-Stage ---
Invitation to Participate}

Hi Team,

I am reaching out to invite you to participate in a research initiative I am
conducting as part of my Master's degree at IME-USP.

\textbf{What is it about?}
We will be exploring a technique called Architectural Hypothesis Engineering
(ArchHypo) and a supporting tool called Hypo-Stage, which I have integrated
into our Backstage portal. The idea is simple: instead of leaving
architectural uncertainties as tacit knowledge held by a few engineers, we
write them down as structured hypotheses, assess their urgency, and track
the actions we are taking to resolve them.

\textbf{What does participation involve?}
I will run a short training session (approximately 60--90 minutes) to introduce
the technique and the tool. After that, you are invited to use Hypo-Stage
whenever you encounter architectural uncertainties in your normal work over
the next four weeks. At the end of the period, I would like to have a brief
conversation (20--30 minutes) with interested team members about their
experience.

\textbf{Is it mandatory?}
No. Participation is entirely voluntary. Using the tool as much or as little
as you find valuable is up to you.

\textbf{What happens to the data?}
All records will be anonymized before being reported in my thesis. The names
of our teams, our products, and our systems will not appear in any research
output. Observations are used only for academic research.

Please reply if you have any questions, or let me know if you would like to
discuss the initiative before the training session.

Looking forward to working with you on this!

[Researcher --- role and name anonymized]
\end{quote}

\paragraph{Editorial note on observation length.}
The invitation above references a four-week horizon. In practice, teams continued to use Hypo-Stage voluntarily beyond that window. Hypothesis registration timestamps in Appendix~\ref{app:hypotheses-data} span 11~February~2026 to 12~April~2026 (61~calendar days, approximately nine weeks), corresponding to five consecutive two-week sprint periods from the earliest to the latest registration in the organization's cadence.

%!TeX root=../thesis.tex

% ============================================================
%  Appendix B — Registered Hypotheses Data (Anonymised)
%  Master's Thesis: Evaluating and Refining Hypo-Stage
%  José Gonçalves Lima Neto — IME-USP, 2026
% ============================================================

\chapter{Registered Hypotheses Data}
\label{app:hypotheses-data}

This appendix presents the complete set of thirteen architectural hypotheses
registered by the two participating teams in Hypo-Stage during the observation
period. All entity references, service names, and internal artefact links have
been anonymised in accordance with the ethical agreement described in
Section~\ref{sec:ethical-considerations}. The anonymisation legend is provided
in Table~\ref{tab:entity-legend}. Technical planning entries are reproduced in
full in Section~\ref{sec:app-technical-plans}.

% ─────────────────────────────────────────────────────────────
\section{Anonymisation Legend}
\label{sec:app-legend}
% ─────────────────────────────────────────────────────────────

Hypotheses in Hypo-Stage are linked to components in the organization's Backstage
catalog through \texttt{entityRef} strings (human-readable labels that identify a
service, worker, platform capability, or template entry). To protect confidentiality
under the ethical agreement summarised in Section~\ref{sec:ethical-considerations},
this appendix does not reproduce real catalog identifiers. Instead, each hypothesis
table in Section~\ref{sec:app-hypothesis-registers} lists \emph{anonymised} entity
references in its \textbf{Entities} row; Table~\ref{tab:entity-legend} explains what
each anonymised string denotes at a functional level, without recovering the original
catalog paths or service names.

The table is intentionally typeset here as \emph{non-floating} material so it remains
with this section even though it occupies most of a page: the goal is to keep the
legend immediately available when reading the hypothesis registers, rather than
deferring it past the large tables in Section~\ref{sec:app-hypothesis-registers}.

\noindent\begin{minipage}{\linewidth}
\centering\small
\setlength{\tabcolsep}{5pt}
\begin{tabular}{p{5.2cm}p{7.8cm}}
\toprule
\textbf{Anonymised entity reference} & \textbf{Description} \\
\midrule
\texttt{catalog: Kotlin + JOOQ + Spring service template}
  & Internal catalog entry for the JVM service scaffold using Kotlin, JOOQ, and
    Spring — intended to replace a legacy functional-language stack in the
    regulated financial domain. \\[4pt]
\texttt{internal-framework: JVM platform stack with Kafka integration}
  & The organization's shared application framework for JVM services, covering
    Kafka exposure conventions, request logging, and platform integrations. \\[4pt]
\texttt{service: secure blob storage — API}
  & HTTP API for storing and governing binary documents with security and
    lifecycle rules. \\[4pt]
\texttt{service: secure blob storage — worker}
  & Background worker for the secure blob storage domain (async processing and
    maintenance). \\[4pt]
\texttt{service: confidential records \& credit-bureau integration — API}
  & API for identity, credit-bureau, and confidential document flows, currently
    coupled to a single business vertical. \\[4pt]
\texttt{service: confidential records \& credit-bureau integration — worker}
  & Worker for the same confidential records and bureau integration context. \\[4pt]
\texttt{worker: high-volume billing aggregation}
  & Long-running worker that aggregates large billing datasets. \\[4pt]
\texttt{platform: shared streaming infrastructure}
  & Shared Kafka streaming platform and consumption standards, modelled as a
    catalog component. \\
\bottomrule
\end{tabular}
\captionof{table}{Anonymised Backstage entity references and their descriptions.}
\label{tab:entity-legend}
\end{minipage}
\par\bigskip

% ─────────────────────────────────────────────────────────────
\section{Hypothesis Registers}
\label{sec:app-hypothesis-registers}
% ─────────────────────────────────────────────────────────────

Each hypothesis is typeset as a single register table: fixed metadata rows, a
full-width statement row, then author notes and technical-plan references (see
Section~\ref{sec:app-technical-plans}).

% #1 ID … #8 Uncertainty/Impact; #9 statement. Then call \hypoboxtail{notes}{plans}.
\newcommand{\hypoboxnine}[9]{%
  \par\addvspace{0.55\baselineskip}%
  \noindent\begingroup
  \small
  \setlength{\tabcolsep}{6pt}%
  \begin{tabular}{@{} >{\bfseries\raggedright\arraybackslash}p{3.4cm} >{\raggedright\arraybackslash}p{\dimexpr\textwidth-3.4cm-6\tabcolsep} @{}}
  \toprule
  \multicolumn{2}{@{}l@{}}{\large #1} \\
  \midrule
  Team & #2 \\
  Registered & #3 \\
  Lifecycle status & #4 \\
  Source type & #5 \\
  Uncertainty / Impact & #8 \\
  Quality attributes & #7 \\
  Entities & {\begingroup\ttfamily\small #6\endgroup} \\
  \addlinespace
  \multicolumn{2}{@{}>{\raggedright\arraybackslash}p{\dimexpr\textwidth-4\tabcolsep}@{}}{%
    \textbf{Statement}\par\vspace{0.3\baselineskip}%
    ``#9''\par
  } \\
  \midrule
}
\newcommand{\hypoboxtail}[2]{%
  Notes & #1 \\
  Technical plans & #2 \\
  \bottomrule
  \end{tabular}%
  \endgroup
  \par
}

% ── H01 ─────────────────────────────────────────────────────
\hypoboxnine%
  {H01}%
  {Product}%
  {2026-03-18}%
  {Open}%
  {Requirements}%
  {catalog: Kotlin + JOOQ + Spring service template}%
  {Performance \textbar{} Reliability \textbar{} Maintainability}%
  {Very High / Very High}%
  {The Kotlin + JOOQ + Spring service template is ready to scale for use in
   regulated financial product services, replacing a legacy functional-language
   stack.}%
\hypoboxtail%
  {The organization intends to deprecate the legacy stack in that domain; the
   Kotlin + JOOQ + Spring template is the leading candidate to replace it.}%
  {None registered at end of observation period.}

% ── H02 ─────────────────────────────────────────────────────
\hypoboxnine%
  {H02}%
  {Platform}%
  {2026-02-11}%
  {Open}%
  {Requirements}%
  {internal-framework: JVM platform stack with Kafka integration}%
  {Configurability \textbar{} Flexibility \textbar{} Extensibility \textbar{} Modifiability}%
  {Very Low / Medium}%
  {Kafka should be exposed to application layers for direct use via the internal
   platform framework.}%
\hypoboxtail%
  {Subject on hold; technical plans expected for the following quarter.}%
  {None registered at end of observation period.}

% ── H03 ─────────────────────────────────────────────────────
\hypoboxnine%
  {H03}%
  {Platform}%
  {2026-02-19}%
  {Open}%
  {Requirements}%
  {internal-framework: JVM platform stack with Kafka integration}%
  {Observability \textbar{} Auditability}%
  {Medium / Very High}%
  {Decouple engineering/observability logs from security/compliance auditing by
   evolving database logging toward a unified Request Logs capability on the
   internal platform framework.}%
\hypoboxtail%
  {The platform framework must let applications derive full log context including
   database logs. It is expected to supersede a legacy audit library and satisfy
   organization-wide security audit requirements.}%
  {TP-03-A, TP-03-B, TP-03-C (see Section~\ref{sec:app-technical-plans}).}

% ── H04 ─────────────────────────────────────────────────────
\hypoboxnine%
  {H04}%
  {Product}%
  {2026-03-18}%
  {Open}%
  {Requirements}%
  {service: secure blob storage — API \textbar{} service: secure blob storage — worker}%
  {Auditability \textbar{} Usability}%
  {Medium / Very Low}%
  {The secure blob storage service must retain a reference to the data subject in
   its domain data so all sensitive blobs can be deleted when the user exercises
   their right to erasure under applicable privacy regulation.}%
\hypoboxtail%
  {Concern raised in an internal RFC about fully erasing user sensitive data,
   including binary objects.}%
  {None registered at end of observation period.}

% ── H05 ─────────────────────────────────────────────────────
\hypoboxnine%
  {H05}%
  {Product}%
  {2026-03-18}%
  {In Review}%
  {Solution}%
  {service: confidential records \& credit-bureau integration — API \textbar{} \ldots{} worker}%
  {Interoperability \textbar{} Extensibility \textbar{} Reusability}%
  {High / Very High}%
  {Credit-bureau integrations can be decoupled from vertical-specific domain
   entities into context-agnostic middleware that accepts only each provider's
   minimum required data (e.g., national ID), reusable across domains without
   duplicated cost.}%
\hypoboxtail%
  {Discovery focused on: (1)~dependency matrix --- current vs.\ real dependency
   per integration; (2)~agnostic interface --- contract without domain concepts,
   only minimal identifiers and bureau outputs; (3)~migration plan --- which flows
   can isolate immediately vs.\ requiring deeper business-logic refactors.}%
  {TP-05-A (see Section~\ref{sec:app-technical-plans}).}

% ── H06 ─────────────────────────────────────────────────────
\hypoboxnine%
  {H06}%
  {Product}%
  {2026-02-24}%
  {Open}%
  {Solution}%
  {worker: high-volume billing aggregation}%
  {Performance \textbar{} Reliability \textbar{} Scalability}%
  {Very Low / Very High}%
  {Chunked keyset pagination with isolated repository-level transactions keeps
   memory strictly bounded by chunk size and prevents out-of-memory failures when
   aggregating millions of billing records.}%
\hypoboxtail%
  {Validated with a capacity/volume test at 5$\times$ production scale (45M rows
   total, 1.5M rows/month). Profiling showed JPA streams can cause first-level
   cache growth; chunking scopes transactions ($\approx$10k rows) so the garbage
   collector can reclaim memory (sawtooth heap pattern). Composite index on
   \texttt{(billing\_year\_month, status, id)}.}%
  {TP-06-A (see Section~\ref{sec:app-technical-plans}).}

% ── H07 ─────────────────────────────────────────────────────
\hypoboxnine%
  {H07}%
  {Product}%
  {2026-04-07}%
  {Validated}%
  {Requirements}%
  {service: secure blob storage — API \textbar{} service: secure blob storage — worker}%
  {Scalability \textbar{} Usability \textbar{} Configurability \textbar{} Flexibility}%
  {Very Low / Very High}%
  {Secure blob storage must support transferring blob ownership across service
   domains so blobs can move from one owning service to another.}%
\hypoboxtail%
  {Validated through a technical discovery spike that defined the transfer
   capability and produced a prioritised implementation backlog.}%
  {TP-07-A (see Section~\ref{sec:app-technical-plans}).}

% ── H08 ─────────────────────────────────────────────────────
\hypoboxnine%
  {H08}%
  {Product}%
  {2026-03-18}%
  {Validated}%
  {Requirements}%
  {service: secure blob storage — API}%
  {Security \textbar{} Usability \textbar{} Interoperability \textbar{} Extensibility}%
  {Very Low / Very High}%
  {Secure blob storage must support sharing a blob across multiple clients so the
   same object can be read from different domain services.}%
\hypoboxtail%
  {Example: a customer-facing channel receives documents later routed to support
   tooling; support needs read access in chat history without breaking vault
   invariants.}%
  {TP-08-A (see Section~\ref{sec:app-technical-plans}).}

% ── H09 ─────────────────────────────────────────────────────
\hypoboxnine%
  {H09}%
  {Platform}%
  {2026-04-12}%
  {Open}%
  {Solution}%
  {platform: shared streaming infrastructure}%
  {Auditability \textbar{} Reliability}%
  {Medium / Very High}%
  {Using the Dead-Letter Topic (DLT) only for auditing and diagnosis, and
   reprocessing from the original topic's offset, preserves consistency and
   auditability better than re-enqueueing messages from the DLT.}%
\hypoboxtail%
  {How to falsify: controlled replay from the DLT without mutating the event must
   be proven safer, simpler, and more auditable than replay by offset.}%
  {TP-09-A, TP-09-B (see Section~\ref{sec:app-technical-plans}).}

% ── H10 ─────────────────────────────────────────────────────
\hypoboxnine%
  {H10}%
  {Platform}%
  {2026-04-12}%
  {Open}%
  {Solution}%
  {platform: shared streaming infrastructure}%
  {Interoperability \textbar{} Performance}%
  {Medium / High}%
  {For short-lived transient failures, limited synchronous retry on the consumer
   thread preserves per-partition ordering better and with lower overall complexity
   than asynchronous retry on a separate queue.}%
\hypoboxtail%
  {How to falsify: a per-partition retry queue preserves order with less blocking
   time and acceptable operational cost.}%
  {TP-10-A (see Section~\ref{sec:app-technical-plans}).}

% ── H11 ─────────────────────────────────────────────────────
\hypoboxnine%
  {H11}%
  {Platform}%
  {2026-04-12}%
  {Open}%
  {Solution}%
  {platform: shared streaming infrastructure}%
  {Reliability \textbar{} Modifiability \textbar{} Interoperability}%
  {Medium / High}%
  {Keeping failure state inside the Kafka message flow increases coordination
   complexity across threads and complicates failure handling, compared to
   handling failures in the execution layer and keeping consumption non-blocking.}%
\hypoboxtail%
  {How to falsify: modelling failure state in Kafka shows lower operational
   complexity and equal or better recovery than handling failures in the execution
   layer.}%
  {TP-11-A (see Section~\ref{sec:app-technical-plans}).}

% ── H12 ─────────────────────────────────────────────────────
\hypoboxnine%
  {H12}%
  {Platform}%
  {2026-04-12}%
  {Open}%
  {Solution}%
  {platform: shared streaming infrastructure}%
  {Performance \textbar{} Reliability}%
  {Medium / High}%
  {Classifying errors as recoverable vs.\ non-recoverable before applying retry
   reduces consumer lag and computational cost without worsening the successful
   processing rate.}%
\hypoboxtail%
  {How to falsify: indiscriminate retries improve final success rate without
   meaningful latency/capacity cost.}%
  {TP-12-A, TP-12-B (see Section~\ref{sec:app-technical-plans}).}

% ── H13 ─────────────────────────────────────────────────────
\hypoboxnine%
  {H13}%
  {Platform}%
  {2026-04-12}%
  {Open}%
  {Solution}%
  {platform: shared streaming infrastructure}%
  {Auditability}%
  {Very Low / Very High}%
  {Allowing fixes by mutating the payload before re-enqueueing increases
   consistency violations and makes traceability and auditing harder.}%
\hypoboxtail%
  {How to falsify: a formal correction mechanism fully preserves audit trail,
   causality, and reproducibility of the event.}%
  {TP-13-A (see Section~\ref{sec:app-technical-plans}).}

\bigskip
\section{Technical Plans Catalog}
\label{sec:app-technical-plans}
% ─────────────────────────────────────────────────────────────

The following entries reproduce all technical plans registered in Hypo-Stage
during the observation period. Each entry is identified by a plan code
(e.g., TP-03-A), the parent hypothesis, the action type, the target date, a
description of the planned action, and the expected outcome. Internal artefact
links (tickets, pull requests, design documents) have been redacted.

\vspace{6pt}

%% ─── H03 plans ────────────────────────────────────────────
\paragraph{TP-03-A \textnormal{(H03 — Decouple eng.\ logs from security/compliance audit)}}
\textit{Action type:} Architectural Spike \quad \textit{Target date:} 2026-03-31 \\
\textit{Description:} Discover and implement default mapping of database
transaction IDs into the Request Log structure within the internal framework
so every request and application log entry carries database transaction IDs for
correlation and audit trail without per-application integration. \\
\textit{Expected outcome:} Request logs consistently include database transaction
ID fields; applications using the framework obtain audit-ready context without
custom wiring.

\paragraph{TP-03-B \textnormal{(H03)}}
\textit{Action type:} Modularisation / Flexibility \quad \textit{Target date:} 2026-03-31 \\
\textit{Description:} Extend framework logging to worker modules so critical
database transactions in workers log with the same context (transaction IDs,
correlation) as request paths. \\
\textit{Expected outcome:} Workers emit logs with database transaction IDs
aligned with the Request Log infrastructure; audit coverage for
asynchronous/background work.

\paragraph{TP-03-C \textnormal{(H03)}}
\textit{Action type:} Modularisation / Flexibility \quad \textit{Target date:} 2026-03-31 \\
\textit{Description:} Stream database transaction properties to the same log
buckets and pipeline as Request and Application logs for a single
audit/observability backbone. \\
\textit{Expected outcome:} Database transaction properties follow the same
ingestion and retention path as Request logs; no separate audit-only pipeline.

%% ─── H05 plan ─────────────────────────────────────────────
\paragraph{TP-05-A \textnormal{(H05 — Credit-bureau decoupling seam)}}
\textit{Action type:} Discovery \quad \textit{Target date:} ongoing \\
\textit{Description:} Structured discovery covering: (1)~dependency matrix
(current vs.\ real dependencies per integration); (2)~agnostic interface
proposal (contract free of domain concepts, accepting only minimal provider
identifiers and bureau outputs); (3)~migration sequencing (which flows can
isolate immediately vs.\ requiring deeper business-logic refactors). \\
\textit{Expected outcome:} A prioritised integration decoupling plan and a
draft interface contract.

%% ─── H06 plan ─────────────────────────────────────────────
\paragraph{TP-06-A \textnormal{(H06 — Billing aggregation chunked keyset pagination)}}
\textit{Action type:} Implementation (Other) \quad \textit{Target date:} 2026-03-13 \\
\textit{Description:} Implement the billing aggregation worker with chunked
keyset pagination using a \texttt{LIMIT} loop with
\texttt{id > :lastProcessedId}; avoid broad \texttt{@Transactional} so the
Hibernate first-level cache clears per chunk; add composite index on
\texttt{(billing\_year\_month, status, id)}. \\
\textit{Expected outcome:} Process 300k+ rows/month safely without
out-of-memory failures; stable sawtooth heap profile; healthy database
CPU/I/O via index scans.

%% ─── H07 plan ─────────────────────────────────────────────
\paragraph{TP-07-A \textnormal{(H07 — Blob ownership transfer)}}
\textit{Action type:} Architectural Spike \quad \textit{Target date:} 2026-04-30 \\
\textit{Description:} Technical discovery to define the transfer capability on
the blob storage API. \\
\textit{Expected outcome:} A backlog for implementing the ownership transfer
feature.

%% ─── H08 plan ─────────────────────────────────────────────
\paragraph{TP-08-A \textnormal{(H08 — Blob cross-client sharing)}}
\textit{Action type:} Guideline Implementation \quad \textit{Target date:} 2026-04-06 \\
\textit{Description:} Alignment session on cross-domain sharing and transfers
so that long-lived access (e.g., messaging history) does not require definitive
ownership by every consumer. \\
\textit{Expected outcome:} A governance decision enabling prioritisation of
sharing and transfer implementation work.

%% ─── H09 plans ────────────────────────────────────────────
\paragraph{TP-09-A \textnormal{(H09 — DLT for audit vs.\ offset replay)}}
\textit{Action type:} Guideline Implementation \quad \textit{Target date:} 2026-06-30 \\
\textit{Description:} Produce a runbook for reprocessing by offset with minimal
tooling for controlled replay. \\
\textit{Expected outcome:} Documentation of how to reprocess by offset with
minimal tooling.

\paragraph{TP-09-B \textnormal{(H09)}}
\textit{Action type:} Experiment \quad \textit{Target date:} 2026-06-30 \\
\textit{Description:} Operational exercise with a simulated incident. \\
\textit{Expected outcome:} Measure MTTR, duplicate rate, and divergence between
the original event and the final state.

%% ─── H10 plan ─────────────────────────────────────────────
\paragraph{TP-10-A \textnormal{(H10 — Synchronous retry for transient failures)}}
\textit{Action type:} Tracer Bullet \quad \textit{Target date:} 2026-06-30 \\
\textit{Description:} Tracer bullet on a critical consumer path; inject
transient failures with windows of 1s, 5s, and 30s. \\
\textit{Expected outcome:} Measure out-of-order delivery, consumer lag,
throughput, and retry success rate.

%% ─── H11 plan ─────────────────────────────────────────────
\paragraph{TP-11-A \textnormal{(H11 — Failure state inside vs.\ outside Kafka flow)}}
\textit{Action type:} Architectural Spike \quad \textit{Target date:} 2026-06-30 \\
\textit{Description:} Spike comparing both failure-state approaches on a real
consumer. \\
\textit{Expected outcome:} Measure consumer lag, blocked-thread time,
states/transitions, and operational recovery steps for each approach.

%% ─── H12 plans ────────────────────────────────────────────
\paragraph{TP-12-A \textnormal{(H12 — Error classification before retry)}}
\textit{Action type:} Guideline Implementation \quad \textit{Target date:} 2026-06-30 \\
\textit{Description:} Define an error taxonomy and mandatory classification step
before retry is applied. \\
\textit{Expected outcome:} Practical directives for classifying errors before
retry, consumable by all teams using the shared streaming platform.

\paragraph{TP-12-B \textnormal{(H12)}}
\textit{Action type:} Software Analytics (Architectural Trigger) \quad \textit{Target date:} 2026-06-30 \\
\textit{Description:} Instrument analytics by error class on a real consumer:
success rate, average attempts, extra time per message, and discards to DLT. \\
\textit{Expected outcome:} A metrics dataset sufficient to falsify or support
the hypothesis.

%% ─── H13 plan ─────────────────────────────────────────────
\paragraph{TP-13-A \textnormal{(H13 — Payload immutability before re-enqueueing)}}
\textit{Action type:} Experiment \quad \textit{Target date:} 2026-06-30 \\
\textit{Description:} Enforce an immutable-payload guideline; add automated
validation in replay tooling; require mandatory architectural review for
exceptions. \\
\textit{Expected outcome:} Experiment results that validate or refute the
hypothesis and establish a baseline for payload governance in the shared
streaming platform.



%%%%%%%%%%%%%%%%%%%%%%%%%%%% ANNEXES %%%%%%%%%%%%%%%%%%%%%%%%%%%%%%%%%%%%%%%%%%%

% Annexes = external documents relevant to the thesis (specs, regulations, etc.)
% Uncomment when needed:
%\cleardoublepage
%\pagestyle{appendix}
%\annex
%\addappheadtotoc
%
% \input{chapters/annex-a-external-docs}


%%%%%%%%%%%%%%%%%%%%%%% BACK MATTER (BIBLIOGRAPHY) %%%%%%%%%%%%%%%%%%%%%%%%%%%%%

\backmatter
\pagestyle{backmatter}

\addtocontents{toc}{\vspace{2\baselineskip plus .5\baselineskip minus .5\baselineskip}}

\printbibliography[
  title=\refname\label{sec:bib},
  heading=bibintoc,
]

\printindex

\end{document}
